% =====================================================================
%  TP2 — Presentación de defensa (10 min + 10 min de preguntas)
%  72.75 Aprendizaje Automático — ITBA — 2026 Q2 — Grupo 7
%
%  Compilar desde informe/, con pdflatex (el motor del TP1), dejando los auxiliares aparte:
%    latexmk -pdf -outdir=_build presentacion.tex && cp _build/presentacion.pdf .
%  El .log se lee antes de borrar _build/.
%
%  PRINCIPIO RECTOR, el del TP1: la slide apoya lo que se dice, no lo dice. Cada frame lleva un
%  \note{} con lo que se dice mientras está en pantalla y, al final, la reserva «Si preguntan».
%  En pantalla quedan la figura, el diagrama o el número, y la frase que cierra la sección. Con
%  las notas ocultas (lo que hace beamer por defecto) el PDF es el de proyección; para ensayar
%  con el guion al lado, descomentar una de las dos líneas \setbeameroption de abajo.
%
%  NÚMEROS (gate G3). Ninguno se escribe a mano, tampoco en las notas: los de train y de
%  validación salen de numeros.tex (src/numeros.py) y los de test, de resultados-test.tex
%  (src/evaluar_test.py). Lo comprueba, desde la raíz del repo:
%    python3 -m src.numeros --verificar informe/presentacion.tex
%  Los macros de los dos archivos llevan la coma entre llaves y los miles con espacio fino, y se
%  pueden usar dentro de $…$, salvo los intervalos de test (\aucic, \recallic): su «--» es una
%  raya sólo en modo texto.
%
%  ESTADO DEL TEST. El test todavía no se evaluó (N0-1): resultados-test.tex declara
%  \testpendientetrue y sus macros valen «?». Las dos slides que dependen del test (14 y 15)
%  tienen escritas sus dos variantes con \iftestpendiente. Correr `python3 -m src.evaluar_test`,
%  después `python3 -m src.numeros`, y recompilar es el único paso para pasar de una a otra.
%
%  ESTRUCTURA. Los 20 frames del esqueleto del plan (§5, ola 7), que suman 9:15, sin slides de
%  respaldo: las preguntas se responden sólo con estas 20 (N0-15). Todo lo calificable,
%  limitaciones incluidas, va antes del hallazgo, y el hallazgo es lo último que se muestra
%  (guía A4, A6).
%
%  Lo que no se trajo del preámbulo del TP1, por código muerto (guía TP1-L): \piefig,
%  \cierreconnumero y \cierrependiente, el interruptor \testevaluado (el TP2 usa el
%  \iftestpendiente de resultados-test.tex), los bloques de beamer, la configuración de tablas
%  de booktabs (el deck no tiene tablas: la matriz de confusión es un dibujo), \chippendiente y
%  la rama del apéndice de la barra, del índice y del pie, que quedó sin uso al sacar el
%  respaldo (N0-15).
% =====================================================================
\documentclass[aspectratio=169,11pt]{beamer}

\usepackage[T1]{fontenc}
\usepackage[utf8]{inputenc}
\usepackage[spanish,es-noquoting]{babel}
% Con babel en español, \% agrega su propio espacio fino, y los macros de numeros.tex ya traen el
% suyo (11{,}3\,\%): sin esta línea el espacio queda doble.
\spanishplainpercent

% ---------------------------------------------------------------------
% TIPOGRAFÍA — Fira Sans + Fira Mono + matemática sans (las razones, en el deck del TP1):
% trazo parejo que sobrevive a la proyección, un solo sistema para el texto y los nombres de
% variable en \texttt, y pesos para armar la jerarquía. `scaled=0.96` devuelve la mancha de
% texto a la de Latin Modern; `lining,tabular` alinea las cifras por dígito.
% ---------------------------------------------------------------------
\usepackage[sfdefault,lining,tabular,scaled=0.96]{FiraSans}
\usepackage[scaled=0.86]{FiraMono}
\usepackage{amsmath}
\usepackage{newtxsf}
\usefonttheme{professionalfonts}
\usepackage[tracking=true]{microtype}

\usepackage{graphicx}
\usepackage{tikz}
\usetikzlibrary{arrows.meta,positioning,fit,backgrounds}

% Sólo las figuras de proyección (guía B5): ninguna figura del deck es la de análisis.
\graphicspath{{../figuras/presentacion/}}

% ---------------------------------------------------------------------
% Tema mínimo, construido a mano: la pantalla es para los datos. Sin símbolos de navegación ni
% transiciones (guía B7).
% ---------------------------------------------------------------------
\usetheme{default}
\usecolortheme{default}
\setbeamertemplate{navigation symbols}{}

% GUION HABLADO. Para ensayar con las notas a la vista:
%   - una página de nota después de cada slide:
% \setbeameroption{show notes}
%   - slide y nota lado a lado:
% \setbeameroption{show notes on second screen=right}

\setbeamersize{text margin left=1.6em,text margin right=1.6em}

% Paleta: la de las figuras (src/estilo.py), validada para daltonismo (tests/test_paleta.py).
% En las figuras, azul es «no», train y barajado; naranja es «yes», validación y hacia adelante.
% En el texto, azul es el dato (\dato) y naranja lo que hay que mirar (\ojo).
\definecolor{azul}{HTML}{2A78D6}
\definecolor{naranja}{HTML}{EB6834}
\definecolor{azuloscuro}{HTML}{104281}
\definecolor{tinta}{HTML}{0B0B0B}
\definecolor{tintasec}{HTML}{52514E}
\definecolor{apagado}{HTML}{898781}
\definecolor{rejilla}{HTML}{E1E0D9}
\definecolor{superficie}{HTML}{FCFCFB}
% El color de cada modelo, como COLOR_MODELO de src/estilo.py.
\colorlet{colorrf}{naranja}
\colorlet{colorknn}{azul}
\definecolor{colornb}{HTML}{A8447F}
\colorlet{colorsvm}{azuloscuro}

\setbeamercolor{normal text}{fg=tinta,bg=superficie}
\setbeamercolor{structure}{fg=azul}
\setbeamercolor{frametitle}{fg=tinta,bg=superficie}
\setbeamercolor{title}{fg=tinta}
\setbeamercolor{alerted text}{fg=naranja}
\setbeamerfont{frametitle}{size=\large,series=\bfseries}
\setbeamerfont{title}{size=\LARGE,series=\bfseries}

% Viñeta: un cuadrado pequeño en gris, el vocabulario rectangular del resto de la pantalla; el
% azul queda para el dato.
\setbeamertemplate{itemize item}{%
  \raisebox{0.18em}{\textcolor{apagado}{\rule{0.38em}{0.38em}}}}

% La mayúscula espaciada, en un solo lugar: rotula un dato (\etiqueta, \cab) o abre una pantalla
% (\volanta, en la portada y en el cierre).
\newcommand{\etiqueta}[1]{\textls[90]{\MakeUppercase{#1}}}
\newcommand{\volanta}[1]{\color{azul}\textls[160]{\MakeUppercase{#1}}}
\newcommand{\cab}[1]{{\footnotesize\color{tintasec}\etiqueta{#1}}}

% ---------------------------------------------------------------------
% ENCABEZADO: título, índice de secciones y barra de progreso (la del TP1).
%
% La barra cuenta frames y no overlays, y divide por \insertmainframenumber. El tope es para la
% primera pasada, cuando todavía no existe el .nav y \insertmainframenumber vale 1.
% ---------------------------------------------------------------------
\newlength{\barraancho}
\newlength{\barrallena}
\newcommand{\barraprogreso}{%
  \setlength{\barraancho}{\textwidth}%
  \pgfmathsetlength{\barrallena}{\barraancho*\insertframenumber/\insertmainframenumber}%
  \ifdim\barrallena>\barraancho\setlength{\barrallena}{\barraancho}\fi
  \textcolor{azul}{\rule{\barrallena}{1.4pt}}%
  \ifdim\barrallena<\barraancho
    \textcolor{rejilla}{\rule{\dimexpr\barraancho-\barrallena\relax}{1.4pt}}%
  \fi
}

% El número de slide, para las preguntas («vuelvan a la 12»): n/N.
\newcommand{\indicediapo}{\insertframenumber/\insertmainframenumber}

% El índice de secciones bajo el título: la sección en curso en azul y negrita, las demás en
% gris. Cada nombre va en una caja del ancho de su versión en negrita, así las posiciones no se
% mueven de una slide a otra.
\newcommand{\sep}{\hspace{0.42em}\textcolor{apagado!40}{\textperiodcentered}\hspace{0.42em}}
\newlength{\secancho}
\newcommand{\secitem}[2]{%
  \settowidth{\secancho}{\textbf{#2}}%
  \makebox[\secancho][c]{%
    \ifnum\insertsectionnumber=#1\relax
      \textcolor{azul}{\textbf{#2}}%
    \else
      \textcolor{apagado}{#2}%
    \fi}}
\newcommand{\indicesecciones}{{\fontsize{6.2}{7}\selectfont%
  \secitem{1}{El problema}\sep\secitem{2}{Datos}\sep\secitem{3}{EDA}\sep
  \secitem{4}{Métricas}\sep\secitem{5}{Modelos}\sep\secitem{6}{Hiperparámetros}\sep
  \secitem{7}{Modelo final}\sep\secitem{8}{Conclusiones}\sep\secitem{9}{Hallazgo}}}

% El encabezado en un \vbox sin \baselineskip entre cajas: su alto es exactamente la suma de los
% \vskip y de las tres cajas. El título va en un \parbox para que uno largo corte en dos líneas
% en lugar de desbordar.
\setbeamertemplate{frametitle}{%
  \vbox{\offinterlineskip%
    \vskip0.55em%
    \parbox[t]{\textwidth}{\raggedright\insertframetitle}%
    \vskip0.60em%
    \hbox{\indicesecciones}%
    \vskip0.32em%
    \hbox{\barraprogreso}%
    \vskip0.15em%
  }%
}

% Pie: la sección a la izquierda, en versalita espaciada, y el número de slide a la derecha, con
% el mismo margen que el texto. El nombre sale de \secname y no de \insertsectionhead, que es un
% \hyperlink: \MakeUppercase le pasaba a mayúsculas también el destino («NAVIGATION2» en lugar
% de «Navigation2») y pdfTeX avisaba un enlace roto por sección (el deck del TP1 tenía el mismo
% aviso).
\setbeamertemplate{footline}{%
  \hbox{\begin{beamercolorbox}[wd=\paperwidth,ht=2.6ex,dp=1.5ex]{}%
    \hspace*{1.6em}%
    {\fontsize{6.2}{7}\selectfont\color{apagado}\etiqueta{\secname}}\hfill
    {\fontsize{6.6}{7}\selectfont\color{apagado}\indicediapo}%
    \hspace*{1.6em}%
  \end{beamercolorbox}}%
}

% ---------------------------------------------------------------------
% FIGURAS. La figura entra en una caja de ancho y alto con keepaspectratio: ata la cota que
% llegue primero. Medido en este deck: bajo un título de una línea quedan unos 193 pt de alto
% útil (\textheight vale 252 pt porque el tema no reserva el encabezado ni el pie), y una figura
% de proyección (2,12:1) a 0,97 del ancho mide 191 pt. Entra sola; con una línea debajo, no. Por
% eso el alto se baja dentro del frame con \renewcommand (el frame es un grupo, no se escapa):
% 0,68 con un veredicto debajo y 0,72 con una referencia. Sin pie (guía B1).
% ---------------------------------------------------------------------
\newcommand{\figgrandealto}{0.90}
\newcommand{\figgrande}[2][0.99]{%
  \includegraphics[width=#1\textwidth,height=\figgrandealto\textheight,keepaspectratio]{#2}}

% Centrado sin el entorno `center`, que agrega \topsep arriba y abajo.
\newcommand{\centrado}[1]{{\par\centering#1\par}}

% Atajos de énfasis.
\newcommand{\dato}[1]{\textcolor{azul}{\textbf{#1}}}
\newcommand{\ojo}[1]{\textcolor{naranja}{\textbf{#1}}}
\newcommand{\gris}[1]{\textcolor{apagado}{#1}}
\newcommand{\sepdato}{\hspace{0.55em}\textcolor{apagado!55}{\textperiodcentered}\hspace{0.55em}}
\newcommand{\salto}{{\color{apagado}$\;\to\;$}}
% En las notas, la marca de avance del guion: una pulsación, un overlay.
\newcommand{\avanza}{\textcolor{naranja}{\,[$\rightarrow$]\,}}

% ---------------------------------------------------------------------
% LA FICHA: fondo neutro y un filete azul a la izquierda, a todo el alto (el trazo de la barra
% de progreso). Es la pieza plana del TP1, dibujada con TikZ para que el filete acompañe a una
% ficha de dos líneas.
%
% EL VEREDICTO es una ficha centrada, del ancho de su texto: la frase que cierra una sección y
% nombra un efecto (guía D1, gate G12). Entra como último overlay de su slide.
%
% LA REFERENCIA es la línea de base en las slides cuya figura no la dibuja porque «sin modelo»
% (AUC 0,5) queda debajo de su eje: el gate G11 pide el rótulo en toda slide con métricas de
% validación. Va pequeña, gris y a la derecha, para que no compita con la figura.
% ---------------------------------------------------------------------
\tikzset{
  filete/.style={path picture={\fill[#1] (path picture bounding box.north west)
    rectangle ([xshift=2.4pt]path picture bounding box.south west);}},
  ficha/.style={fill=rejilla!45, filete=azul, inner xsep=0.75em, inner ysep=0.5em,
    text width=#1, align=flush left, font=\small\bfseries, text=azuloscuro},
}
\newlength{\fichaancho}
\newcommand{\ficha}[2][\linewidth]{%
  \setlength{\fichaancho}{#1}\addtolength{\fichaancho}{-1.5em}%
  \tikz[baseline]\node[ficha=\fichaancho]{#2};}
\newcommand{\veredicto}[1]{%
  \centrado{\tikz\node[fill=rejilla!45, filete=azul, inner xsep=0.85em, inner ysep=0.45em,
    font=\small\bfseries, text=tinta]{#1};}}
\newcommand{\referencia}[1]{{\par\raggedleft\scriptsize\color{apagado}#1\par}}

% ---------------------------------------------------------------------
% EL NÚMERO HÉROE (slide 14), a 40 pt con su pie: \heroe cuando el test ya se evaluó, y
% \heropendiente, el hueco dibujado como tal (una caja de trazo discontinuo con el rótulo
% adentro), mientras no. El pie de \heropendiente va pegado a su caja: la caja ya delimita.
% ---------------------------------------------------------------------
\newcommand{\heroe}[2]{%
  \begin{center}
    {\fontsize{40}{44}\selectfont\textcolor{azul}{\textbf{#1}}}\\[0.55em]
    {\color{apagado}\rule{2.2em}{1pt}}\\[0.6em]
    {\footnotesize\color{tintasec}#2}
  \end{center}%
}
\newcommand{\heropendiente}[2]{%
  \centrado{%
    \begin{tikzpicture}
      \node[draw=apagado, dash pattern=on 3.2pt off 2.6pt, line width=0.9pt,
            rounded corners=3pt, fill=rejilla!20, inner xsep=1.5em, inner ysep=0.7em]
        {\fontsize{21}{25}\selectfont\color{apagado}\textls[140]{\MakeUppercase{#1}}};
    \end{tikzpicture}\\[0.45em]
    {\footnotesize\color{tintasec}#2}}}

% ---------------------------------------------------------------------
% DIAGRAMAS: un solo juego de estilos para todos (caja neutra, azul y naranja, misma geometría
% y misma flecha). `visible on` es el truco del TP1: el nodo existe en todos los overlays y sólo
% se pinta transparente cuando no toca, así el posicionamiento no falla y el dibujo no salta.
% ---------------------------------------------------------------------
\tikzset{
  caja/.style       = {draw=apagado, fill=rejilla!30, rounded corners=2pt, line width=0.7pt,
                       text=tinta, align=center, inner sep=5pt},
  cajaazul/.style   = {caja, draw=azul, fill=azul!10},
  cajanaranja/.style= {caja, draw=naranja, fill=naranja!10},
  flecha/.style     = {-{Stealth[length=2.2mm,width=1.6mm]}, draw=apagado, line width=0.7pt},
  invisible/.style  = {opacity=0, text opacity=0},
  visible on/.style = {alt={#1{}{invisible}}},
  alt/.code args    = {<#1>#2#3}{\alt<#1>{\pgfkeysalso{#2}}{\pgfkeysalso{#3}}},
}

\title{¿Quién va a contratar un plazo fijo?}
% Sin \textperiodcentered: dentro de \textls (la volanta) deja «minus0pt» impreso en la
% fuente TS1.
\subtitle{Trabajo Práctico 2: clasificación supervisada}
\author{Grupo 7 --- Andrés Cortese (64612) \textperiodcentered{} Sebastián Caules (64331)}
\institute{Instituto Tecnológico de Buenos Aires --- 72.75 Aprendizaje Automático}
\date{Defensa: 14/10/2026}

% Metadatos y apertura: el visor abre sólo la página, con la slide entera (las razones, en el
% deck del TP1).
\hypersetup{
  pdfauthor={Grupo 7 --- Andrés Cortese (64612), Sebastián Caules (64331)},
  pdfsubject={TP2 --- Clasificación supervisada sobre Bank Marketing: Naive Bayes, SVM, KNN y
              Random Forest, con validación cruzada},
  pdfkeywords={clasificación, Naive Bayes, SVM, KNN, Random Forest, validación cruzada, AUC,
               recall, fuga de datos, validación temporal, ITBA, 72.75},
  pdfpagemode=UseNone,
  pdfstartview=Fit,
}

% Los números: todos generados. numeros.tex define además \signoMenos.
% GENERADO POR src/numeros.py; no editar a mano. Se regenera con: python3 -m src.numeros
% Cada macro es una cifra de resultados/ (o de train, para el EDA) en el formato del deck:
% coma decimal, miles y % con espacio fino, menos tipográfico. Su fuente está en el
% comentario de su línea y, legible, en informe/numeros.md. Los números de test están en
% informe/resultados-test.tex (src/evaluar_test.py).
% Se leen igual en texto y dentro de $…$: la coma decimal va como {,} (en modo matemático
% una coma suelta deja un espacio detrás) y el menos es \signoMenos, definido aquí: en
% texto es \textminus, que dentro de $…$ LaTeX compone mal, y en matemática es «-».
% Con babel-spanish, \% ya agrega su propio espacio fino y \,\% queda doble: el preámbulo
% del deck lleva \spanishplainpercent después de \usepackage[spanish]{babel}.
% Valen «?» porque falta su entrada: resultados/evaluacion_test.json (4).

\DeclareRobustCommand{\signoMenos}{\ifmmode-\else\textminus\fi}

% --- Slide 2 · El problema -------------------------------------------------------------------------
\newcommand{\filasCsv}{41\,188}                             % data/raw/bank-additional-names.txt, §5: Number of Instances
\newcommand{\predictoras}{20}                               % data/raw/bank-additional-names.txt, §6: Number of Attributes (sin y)
\newcommand{\filasTrain}{32\,940}                           % train (cargar_train()): filas
\newcommand{\yesTrain}{3\,711}                              % train (cargar_train()): filas con y = yes
\newcommand{\noTrain}{29\,229}                              % train (cargar_train()): filas con y = no
\newcommand{\pctYesTrain}{11{,}3\,\%}                       % train (cargar_train()): % de y = yes
\newcommand{\exactitudSiempreNo}{88{,}7\,\%}                % train (cargar_train()): % de y = no, la exactitud de decir siempre «no»
\newcommand{\razonNoYes}{7{,}9}                             % train (cargar_train()): filas «no» por cada «yes»

% --- Slide 3 · Qué datos deciden -------------------------------------------------------------------
\newcommand{\filasSinDuplicados}{41\,176}                   % src/evaluar_test.py: N_TOTAL, el CSV sin duplicados (D-01)
\newcommand{\duplicados}{12}                                % data/raw/bank-additional-names.txt, §5, menos N_TOTAL de src/evaluar_test.py (D-01)
\newcommand{\filasTest}{8\,236}                             % src/evaluar_test.py: N_TOTAL menos filasTrain (D-02)
\newcommand{\llamadasTest}{1\,647}                          % src/metricas.py: llamadas(filasTest), el presupuesto de D-20
\newcommand{\proporcionTrain}{80\,\%}                       % src/datos.py: 1 − PROP_TEST
\newcommand{\proporcionTest}{20\,\%}                        % src/datos.py: PROP_TEST
\newcommand{\kFolds}{5}                                     % resultados/modelo_elegido.json: k_folds
\newcommand{\semilla}{42}                                   % resultados/modelo_elegido.json: semilla
\newcommand{\filasFoldEntrenamiento}{26\,352}               % resultados/robustez_folds.csv: barajado, fold 1, n_train
\newcommand{\filasFoldValidacion}{6\,588}                   % resultados/robustez_folds.csv: barajado, fold 1, n_validacion
\newcommand{\llamadasFoldValidacion}{1\,318}                % src/metricas.py: llamadas(filasFoldValidacion)
\newcommand{\llamadasTrain}{6\,588}                         % src/metricas.py: llamadas(filasTrain), el corte de la lista fuera de fold
\newcommand{\semillasSinEstratificar}{400}                  % resultados/evidencia_particion.json: n_semillas, las particiones sin estratificar (D-02)
\newcommand{\desvioPctYesTestSinEstratificar}{0{,}31}       % resultados/evidencia_particion.json: sin_estratificar.desvio_pp, el desvío del % de «yes» del test entre semillas, en puntos porcentuales (D-02)
\newcommand{\pctYesTestSinEstratificarMinimo}{10{,}4\,\%}   % resultados/evidencia_particion.json: sin_estratificar.minimo_pct, el menor % de «yes» del test entre semillas (D-02)
\newcommand{\pctYesTestSinEstratificarMaximo}{12{,}1\,\%}   % resultados/evidencia_particion.json: sin_estratificar.maximo_pct, el mayor (D-02)
\newcommand{\pctYesTrainTemporal}{6{,}4\,\%}                % resultados/evidencia_particion.json: temporal.pct_yes_train, % de «yes» de train si el test fuera el último 20 % del archivo (D-03)
\newcommand{\pctYesTestTemporal}{30{,}8\,\%}                % resultados/evidencia_particion.json: temporal.pct_yes_test, el del test (D-03)
\newcommand{\vectoresRepetidos}{1\,153}                     % train (cargar_train()): vectores de las predictoras disponibles antes de llamar (todas las columnas salvo fila, y y duration; D-07) que aparecen en dos filas o más
\newcommand{\filasVectoresRepetidos}{2\,482}                % train (cargar_train()): vectores de las predictoras disponibles antes de llamar (todas las columnas salvo fila, y y duration; D-07): filas de los vectores repetidos
\newcommand{\pctFilasVectoresRepetidos}{7{,}5\,\%}          % train (cargar_train()): vectores de las predictoras disponibles antes de llamar (todas las columnas salvo fila, y y duration; D-07): ídem, en % de las filas de train
\newcommand{\vectoresRepetidosClasesDistintas}{156}         % train (cargar_train()): vectores de las predictoras disponibles antes de llamar (todas las columnas salvo fila, y y duration; D-07): repetidos con «yes» y «no» a la vez
\newcommand{\filasVectoresRepetidosClasesDistintas}{350}    % train (cargar_train()): vectores de las predictoras disponibles antes de llamar (todas las columnas salvo fila, y y duration; D-07): filas de esos vectores
\newcommand{\vectoresRepetidosConDuration}{0}               % train (cargar_train()): vectores repetidos con todas las predictoras, duration incluida (tras quitar los duplicados exactos, D-01)

% --- Slide 4 · EDA: la duración es fuga ------------------------------------------------------------
\newcommand{\durationDecilCorto}{59}                        % train (cargar_train()): duration, borde superior del primer decil (s)
\newcommand{\durationDecilLargo}{551}                       % train (cargar_train()): duration, borde inferior del último decil (s)
\newcommand{\pctYesDurationDecilCorto}{0{,}0\,\%}           % train (cargar_train()): % de «yes» en el primer decil de duration
\newcommand{\pctYesDurationDecilLargo}{46{,}1\,\%}          % train (cargar_train()): % de «yes» en el último decil de duration
\newcommand{\filasDurationCero}{4}                          % train (cargar_train()): filas con duration = 0 (todas «no»)
\newcommand{\aucDurationSola}{0{,}818}                      % train (cargar_train()): AUC de duration sola como puntaje
\newcommand{\predictorasDisponibles}{19}                    % data/raw/bank-additional-names.txt, §6, menos las predictoras que src/datos.py deja fuera del modelo (EXCLUIDAS sin fila: duration); D-05 y D-07
\newcommand{\aucRfConDuration}{0{,}939}                     % resultados/ablaciones_resumen.csv: rf, A1, auc_variante (techo con duration)
\newcommand{\aucNbGaussianoConDuration}{0{,}831}            % resultados/ablaciones_resumen.csv: nb_gaussiano, A1, auc_variante (techo con duration)
\newcommand{\aucSvmConDuration}{0{,}907}                    % resultados/ablaciones_resumen.csv: svm, A1, auc_variante (techo con duration)
\newcommand{\aucKnnConDuration}{0{,}911}                    % resultados/ablaciones_resumen.csv: knn, A1, auc_variante (techo con duration)
\newcommand{\deltaAucDurationNbGaussiano}{+0{,}063}         % resultados/ablaciones_resumen.csv: nb_gaussiano, A1, delta_media (ΔAUC pareado contra A0)
\newcommand{\deltaAucDurationNbGaussianoDesvio}{0{,}0055}   % resultados/ablaciones_resumen.csv: nb_gaussiano, A1, delta_desvio
\newcommand{\deltaAucDurationSvm}{+0{,}206}                 % resultados/ablaciones_resumen.csv: svm, A1, delta_media (ΔAUC pareado contra A0)
\newcommand{\deltaAucDurationSvmDesvio}{0{,}016}            % resultados/ablaciones_resumen.csv: svm, A1, delta_desvio
\newcommand{\deltaAucDurationKnn}{+0{,}156}                 % resultados/ablaciones_resumen.csv: knn, A1, delta_media (ΔAUC pareado contra A0)
\newcommand{\deltaAucDurationKnnDesvio}{0{,}013}            % resultados/ablaciones_resumen.csv: knn, A1, delta_desvio
\newcommand{\deltaAucDurationRf}{+0{,}170}                  % resultados/ablaciones_resumen.csv: rf, A1, delta_media (ΔAUC pareado contra A0)
\newcommand{\deltaAucDurationRfDesvio}{0{,}0072}            % resultados/ablaciones_resumen.csv: rf, A1, delta_desvio

% --- Slide 5 · EDA: el archivo está ordenado por fecha ---------------------------------------------
\newcommand{\anioInicio}{2008}                              % data/raw/bank-additional-names.txt, §4: primer año del archivo
\newcommand{\anioFin}{2010}                                 % data/raw/bank-additional-names.txt, §4: último año del archivo
\newcommand{\filasBloqueEda}{2\,000}                        % src/eda_html.py: BLOQUE, filas del CSV original por bloque
\newcommand{\pctYesPrimerBloqueEda}{1{,}9\,\%}              % train (cargar_train()): % de «yes» en el primer bloque de BLOQUE filas
\newcommand{\pctYesUltimoBloqueEda}{50{,}9\,\%}             % train (cargar_train()): % de «yes» en el último bloque de BLOQUE filas
\newcommand{\factorYesBloquesEda}{27{,}1}                   % train (cargar_train()): % de «yes» del último bloque / del primero
\newcommand{\pctYesDosMilOcho}{4{,}9\,\%}                   % train (cargar_train()): % de «yes» en 2008 (año inferido del orden de los meses, como eda.html)
\newcommand{\pctYesDosMilNueve}{19{,}2\,\%}                 % train (cargar_train()): % de «yes» en 2009 (año inferido del orden de los meses, como eda.html)
\newcommand{\pctYesDosMilDiez}{51{,}9\,\%}                  % train (cargar_train()): % de «yes» en 2010 (año inferido del orden de los meses, como eda.html)
\newcommand{\factorYesAnios}{10{,}5}                        % train (cargar_train()): % de «yes» de 2010 / de 2008
\newcommand{\euriborPrimerBloqueEda}{4{,}86}                % train (cargar_train()): euribor3m medio del primer bloque de BLOQUE filas
\newcommand{\euriborMinimoBloqueEda}{0{,}71}                % train (cargar_train()): euribor3m medio del bloque más bajo
\newcommand{\euriborFoldSinBarajarPrimero}{4{,}87}          % train (cargar_train()): euribor3m medio del fold 1 con StratifiedKFold(5) sin barajar, cv=5 (D-06)
\newcommand{\euriborFoldSinBarajarUltimo}{1{,}10}           % train (cargar_train()): ídem, fold 5 (D-06)
\newcommand{\euriborFoldBarajadoMinimo}{3{,}58}             % train (cargar_train()): euribor3m medio, el menor de los folds() (D-06)
\newcommand{\euriborFoldBarajadoMaximo}{3{,}64}             % train (cargar_train()): ídem, el mayor (D-06)
\newcommand{\pctSeptiembreUltimoTramo}{100{,}0\,\%}         % train (cargar_train()): % de las filas de month = sep que caen en el último 20 % de train por fecha (A8)
\newcommand{\pctOctubreUltimoTramo}{90{,}9\,\%}             % train (cargar_train()): % de las filas de month = oct que caen en el último 20 % de train por fecha (A8)
\newcommand{\pctDiciembreUltimoTramo}{94{,}5\,\%}           % train (cargar_train()): % de las filas de month = dec que caen en el último 20 % de train por fecha (A8)

% --- Slide 6 · EDA: 999 no es «nunca contactado» ---------------------------------------------------
\newcommand{\pdaysCentinela}{999}                           % src/eda_html.py: PDAYS_CENTINELA, el pdays de quien no fue contactado en una campaña anterior (D-10)
\newcommand{\pctPdaysCentinela}{96{,}3\,\%}                 % train (cargar_train()): % de filas con pdays = 999
\newcommand{\filasCentinelaConPrevio}{3\,302}               % train (cargar_train()): filas con pdays = 999 y previous >= 1
\newcommand{\filasPdaysContactados}{1\,207}                 % train (cargar_train()): filas con pdays < 999
\newcommand{\pctYesPdaysContactados}{64{,}6\,\%}            % train (cargar_train()): % de «yes» con pdays < 999
\newcommand{\pctYesPdaysCentinela}{9{,}2\,\%}               % train (cargar_train()): % de «yes» con pdays = 999
\newcommand{\pctPoutcomeInexistente}{86{,}3\,\%}            % train (cargar_train()): % de filas con poutcome = nonexistent
\newcommand{\pctYesPoutcomeExito}{66{,}0\,\%}               % train (cargar_train()): % de «yes» con poutcome = success
\newcommand{\deltaAucSinPdaysNbGaussiano}{\signoMenos 0{,}0012} % resultados/ablaciones_resumen.csv: nb_gaussiano, A2, delta_media (ΔAUC pareado contra A0)
\newcommand{\deltaAucSinPdaysNbGaussianoDesvio}{0{,}0003}   % resultados/ablaciones_resumen.csv: nb_gaussiano, A2, delta_desvio
\newcommand{\deltaAucSinPdaysSvm}{+0{,}0023}                % resultados/ablaciones_resumen.csv: svm, A2, delta_media (ΔAUC pareado contra A0)
\newcommand{\deltaAucSinPdaysSvmDesvio}{0{,}0037}           % resultados/ablaciones_resumen.csv: svm, A2, delta_desvio
\newcommand{\deltaAucSinPdaysKnn}{\signoMenos 0{,}0019}     % resultados/ablaciones_resumen.csv: knn, A2, delta_media (ΔAUC pareado contra A0)
\newcommand{\deltaAucSinPdaysKnnDesvio}{0{,}0007}           % resultados/ablaciones_resumen.csv: knn, A2, delta_desvio
\newcommand{\deltaAucSinPdaysRf}{+0{,}0003}                 % resultados/ablaciones_resumen.csv: rf, A2, delta_media (ΔAUC pareado contra A0)
\newcommand{\deltaAucSinPdaysRfDesvio}{0{,}0022}            % resultados/ablaciones_resumen.csv: rf, A2, delta_desvio

% --- Slide 7 · EDA: el pipeline y la consecuencia de cada decisión ---------------------------------
\newcommand{\columnasReferencia}{62}                        % src/preproceso.py: columnas del one-hot con Opciones(), la referencia de las ablaciones
\newcommand{\columnasFinales}{58}                           % src/preproceso.py: columnas del one-hot con OPCIONES_FINALES (src/configuracion.py)
\newcommand{\columnasDuration}{63}                          % resultados/ablaciones_resumen.csv: A1, columnas
\newcommand{\columnasSinPdays}{61}                          % resultados/ablaciones_resumen.csv: A2, columnas
\newcommand{\columnasDefaultIndicadora}{60}                 % resultados/ablaciones_resumen.csv: A3, columnas
\newcommand{\columnasRaras}{60}                             % resultados/ablaciones_resumen.csv: A4, columnas
\newcommand{\columnasLogCampaign}{62}                       % resultados/ablaciones_resumen.csv: A5, columnas
\newcommand{\columnasMacroReducido}{59}                     % resultados/ablaciones_resumen.csv: A6, columnas
\newcommand{\columnasSinMacro}{57}                          % resultados/ablaciones_resumen.csv: A7, columnas
\newcommand{\columnasSinMacroNiMonth}{47}                   % resultados/ablaciones_resumen.csv: A8, columnas
\newcommand{\columnasSinDiaSemana}{57}                      % resultados/ablaciones_resumen.csv: A9, columnas
\newcommand{\columnasEdadTramos}{70}                        % resultados/ablaciones_resumen.csv: A10, columnas
\newcommand{\columnasImputarModa}{56}                       % resultados/ablaciones_resumen.csv: A11, columnas
\newcommand{\columnasSinEscalar}{62}                        % resultados/ablaciones_resumen.csv: A12, columnas
\newcommand{\aucReferenciaAblacionRf}{0{,}770}              % resultados/ablaciones_resumen.csv: rf, auc_referencia (A0)
\newcommand{\aucReferenciaAblacionNbGaussiano}{0{,}768}     % resultados/ablaciones_resumen.csv: nb_gaussiano, auc_referencia (A0)
\newcommand{\aucReferenciaAblacionSvm}{0{,}701}             % resultados/ablaciones_resumen.csv: svm, auc_referencia (A0)
\newcommand{\aucReferenciaAblacionKnn}{0{,}755}             % resultados/ablaciones_resumen.csv: knn, auc_referencia (A0)
\newcommand{\aucRfSinMacro}{0{,}742}                        % resultados/ablaciones_resumen.csv: rf, A7, auc_variante
\newcommand{\aucNbGaussianoSinMacro}{0{,}746}               % resultados/ablaciones_resumen.csv: nb_gaussiano, A7, auc_variante
\newcommand{\aucSvmSinMacro}{0{,}697}                       % resultados/ablaciones_resumen.csv: svm, A7, auc_variante
\newcommand{\aucKnnSinMacro}{0{,}708}                       % resultados/ablaciones_resumen.csv: knn, A7, auc_variante
\newcommand{\aucRfSinMacroNiMonth}{0{,}676}                 % resultados/ablaciones_resumen.csv: rf, A8, auc_variante
\newcommand{\aucNbGaussianoSinMacroNiMonth}{0{,}718}        % resultados/ablaciones_resumen.csv: nb_gaussiano, A8, auc_variante
\newcommand{\aucSvmSinMacroNiMonth}{0{,}653}                % resultados/ablaciones_resumen.csv: svm, A8, auc_variante
\newcommand{\aucKnnSinMacroNiMonth}{0{,}679}                % resultados/ablaciones_resumen.csv: knn, A8, auc_variante
\newcommand{\deltaAucDefaultIndicadoraNbGaussiano}{\signoMenos 0{,}0003} % resultados/ablaciones_resumen.csv: nb_gaussiano, A3, delta_media (ΔAUC pareado contra A0)
\newcommand{\deltaAucDefaultIndicadoraNbGaussianoDesvio}{0{,}0017} % resultados/ablaciones_resumen.csv: nb_gaussiano, A3, delta_desvio
\newcommand{\deltaAucDefaultIndicadoraSvm}{+0{,}0004}       % resultados/ablaciones_resumen.csv: svm, A3, delta_media (ΔAUC pareado contra A0)
\newcommand{\deltaAucDefaultIndicadoraSvmDesvio}{0{,}0012}  % resultados/ablaciones_resumen.csv: svm, A3, delta_desvio
\newcommand{\deltaAucDefaultIndicadoraKnn}{+0{,}0000}       % resultados/ablaciones_resumen.csv: knn, A3, delta_media (ΔAUC pareado contra A0)
\newcommand{\deltaAucDefaultIndicadoraKnnDesvio}{0{,}0028}  % resultados/ablaciones_resumen.csv: knn, A3, delta_desvio
\newcommand{\deltaAucDefaultIndicadoraRf}{+0{,}0029}        % resultados/ablaciones_resumen.csv: rf, A3, delta_media (ΔAUC pareado contra A0)
\newcommand{\deltaAucDefaultIndicadoraRfDesvio}{0{,}0021}   % resultados/ablaciones_resumen.csv: rf, A3, delta_desvio
\newcommand{\deltaAucRarasNbGaussiano}{+0{,}0008}           % resultados/ablaciones_resumen.csv: nb_gaussiano, A4, delta_media (ΔAUC pareado contra A0)
\newcommand{\deltaAucRarasNbGaussianoDesvio}{0{,}0006}      % resultados/ablaciones_resumen.csv: nb_gaussiano, A4, delta_desvio
\newcommand{\deltaAucRarasSvm}{+0{,}0002}                   % resultados/ablaciones_resumen.csv: svm, A4, delta_media (ΔAUC pareado contra A0)
\newcommand{\deltaAucRarasSvmDesvio}{0{,}0007}              % resultados/ablaciones_resumen.csv: svm, A4, delta_desvio
\newcommand{\deltaAucRarasKnn}{+0{,}0001}                   % resultados/ablaciones_resumen.csv: knn, A4, delta_media (ΔAUC pareado contra A0)
\newcommand{\deltaAucRarasKnnDesvio}{0{,}0006}              % resultados/ablaciones_resumen.csv: knn, A4, delta_desvio
\newcommand{\deltaAucRarasRf}{\signoMenos 0{,}0000}         % resultados/ablaciones_resumen.csv: rf, A4, delta_media (ΔAUC pareado contra A0)
\newcommand{\deltaAucRarasRfDesvio}{0{,}0014}               % resultados/ablaciones_resumen.csv: rf, A4, delta_desvio
\newcommand{\deltaAucLogCampaignSvm}{\signoMenos 0{,}0006}  % resultados/ablaciones_resumen.csv: svm, A5, delta_media (ΔAUC pareado contra A0)
\newcommand{\deltaAucLogCampaignSvmDesvio}{0{,}0020}        % resultados/ablaciones_resumen.csv: svm, A5, delta_desvio
\newcommand{\deltaAucLogCampaignKnn}{+0{,}0008}             % resultados/ablaciones_resumen.csv: knn, A5, delta_media (ΔAUC pareado contra A0)
\newcommand{\deltaAucLogCampaignKnnDesvio}{0{,}0026}        % resultados/ablaciones_resumen.csv: knn, A5, delta_desvio
\newcommand{\deltaAucMacroReducidoNbGaussiano}{\signoMenos 0{,}0041} % resultados/ablaciones_resumen.csv: nb_gaussiano, A6, delta_media (ΔAUC pareado contra A0)
\newcommand{\deltaAucMacroReducidoNbGaussianoDesvio}{0{,}0016} % resultados/ablaciones_resumen.csv: nb_gaussiano, A6, delta_desvio
\newcommand{\deltaAucMacroReducidoSvm}{+0{,}0030}           % resultados/ablaciones_resumen.csv: svm, A6, delta_media (ΔAUC pareado contra A0)
\newcommand{\deltaAucMacroReducidoSvmDesvio}{0{,}0083}      % resultados/ablaciones_resumen.csv: svm, A6, delta_desvio
\newcommand{\deltaAucMacroReducidoKnn}{\signoMenos 0{,}0083} % resultados/ablaciones_resumen.csv: knn, A6, delta_media (ΔAUC pareado contra A0)
\newcommand{\deltaAucMacroReducidoKnnDesvio}{0{,}0030}      % resultados/ablaciones_resumen.csv: knn, A6, delta_desvio
\newcommand{\deltaAucMacroReducidoRf}{+0{,}0005}            % resultados/ablaciones_resumen.csv: rf, A6, delta_media (ΔAUC pareado contra A0)
\newcommand{\deltaAucMacroReducidoRfDesvio}{0{,}0028}       % resultados/ablaciones_resumen.csv: rf, A6, delta_desvio
\newcommand{\deltaAucSinMacroNbGaussiano}{\signoMenos 0{,}023} % resultados/ablaciones_resumen.csv: nb_gaussiano, A7, delta_media (ΔAUC pareado contra A0)
\newcommand{\deltaAucSinMacroNbGaussianoDesvio}{0{,}0037}   % resultados/ablaciones_resumen.csv: nb_gaussiano, A7, delta_desvio
\newcommand{\deltaAucSinMacroSvm}{\signoMenos 0{,}0037}     % resultados/ablaciones_resumen.csv: svm, A7, delta_media (ΔAUC pareado contra A0)
\newcommand{\deltaAucSinMacroSvmDesvio}{0{,}011}            % resultados/ablaciones_resumen.csv: svm, A7, delta_desvio
\newcommand{\deltaAucSinMacroKnn}{\signoMenos 0{,}048}      % resultados/ablaciones_resumen.csv: knn, A7, delta_media (ΔAUC pareado contra A0)
\newcommand{\deltaAucSinMacroKnnDesvio}{0{,}0080}           % resultados/ablaciones_resumen.csv: knn, A7, delta_desvio
\newcommand{\deltaAucSinMacroRf}{\signoMenos 0{,}028}       % resultados/ablaciones_resumen.csv: rf, A7, delta_media (ΔAUC pareado contra A0)
\newcommand{\deltaAucSinMacroRfDesvio}{0{,}0041}            % resultados/ablaciones_resumen.csv: rf, A7, delta_desvio
\newcommand{\deltaAucSinMacroNiMonthNbGaussiano}{\signoMenos 0{,}050} % resultados/ablaciones_resumen.csv: nb_gaussiano, A8, delta_media (ΔAUC pareado contra A0)
\newcommand{\deltaAucSinMacroNiMonthNbGaussianoDesvio}{0{,}0046} % resultados/ablaciones_resumen.csv: nb_gaussiano, A8, delta_desvio
\newcommand{\deltaAucSinMacroNiMonthSvm}{\signoMenos 0{,}048} % resultados/ablaciones_resumen.csv: svm, A8, delta_media (ΔAUC pareado contra A0)
\newcommand{\deltaAucSinMacroNiMonthSvmDesvio}{0{,}0060}    % resultados/ablaciones_resumen.csv: svm, A8, delta_desvio
\newcommand{\deltaAucSinMacroNiMonthKnn}{\signoMenos 0{,}076} % resultados/ablaciones_resumen.csv: knn, A8, delta_media (ΔAUC pareado contra A0)
\newcommand{\deltaAucSinMacroNiMonthKnnDesvio}{0{,}0073}    % resultados/ablaciones_resumen.csv: knn, A8, delta_desvio
\newcommand{\deltaAucSinMacroNiMonthRf}{\signoMenos 0{,}093} % resultados/ablaciones_resumen.csv: rf, A8, delta_media (ΔAUC pareado contra A0)
\newcommand{\deltaAucSinMacroNiMonthRfDesvio}{0{,}0048}     % resultados/ablaciones_resumen.csv: rf, A8, delta_desvio
\newcommand{\deltaAucSinDiaSemanaNbGaussiano}{\signoMenos 0{,}0000} % resultados/ablaciones_resumen.csv: nb_gaussiano, A9, delta_media (ΔAUC pareado contra A0)
\newcommand{\deltaAucSinDiaSemanaNbGaussianoDesvio}{0{,}0002} % resultados/ablaciones_resumen.csv: nb_gaussiano, A9, delta_desvio
\newcommand{\deltaAucSinDiaSemanaSvm}{\signoMenos 0{,}0074} % resultados/ablaciones_resumen.csv: svm, A9, delta_media (ΔAUC pareado contra A0)
\newcommand{\deltaAucSinDiaSemanaSvmDesvio}{0{,}0021}       % resultados/ablaciones_resumen.csv: svm, A9, delta_desvio
\newcommand{\deltaAucSinDiaSemanaKnn}{\signoMenos 0{,}0014} % resultados/ablaciones_resumen.csv: knn, A9, delta_media (ΔAUC pareado contra A0)
\newcommand{\deltaAucSinDiaSemanaKnnDesvio}{0{,}0061}       % resultados/ablaciones_resumen.csv: knn, A9, delta_desvio
\newcommand{\deltaAucSinDiaSemanaRf}{\signoMenos 0{,}0034}  % resultados/ablaciones_resumen.csv: rf, A9, delta_media (ΔAUC pareado contra A0)
\newcommand{\deltaAucSinDiaSemanaRfDesvio}{0{,}0038}        % resultados/ablaciones_resumen.csv: rf, A9, delta_desvio
\newcommand{\deltaAucEdadTramosNbGaussiano}{+0{,}0001}      % resultados/ablaciones_resumen.csv: nb_gaussiano, A10, delta_media (ΔAUC pareado contra A0)
\newcommand{\deltaAucEdadTramosNbGaussianoDesvio}{0{,}0022} % resultados/ablaciones_resumen.csv: nb_gaussiano, A10, delta_desvio
\newcommand{\deltaAucImputarModaNbGaussiano}{+0{,}0000}     % resultados/ablaciones_resumen.csv: nb_gaussiano, A11, delta_media (ΔAUC pareado contra A0)
\newcommand{\deltaAucImputarModaNbGaussianoDesvio}{0{,}0039} % resultados/ablaciones_resumen.csv: nb_gaussiano, A11, delta_desvio
\newcommand{\deltaAucImputarModaSvm}{\signoMenos 0{,}0002}  % resultados/ablaciones_resumen.csv: svm, A11, delta_media (ΔAUC pareado contra A0)
\newcommand{\deltaAucImputarModaSvmDesvio}{0{,}0049}        % resultados/ablaciones_resumen.csv: svm, A11, delta_desvio
\newcommand{\deltaAucImputarModaKnn}{+0{,}0001}             % resultados/ablaciones_resumen.csv: knn, A11, delta_media (ΔAUC pareado contra A0)
\newcommand{\deltaAucImputarModaKnnDesvio}{0{,}0022}        % resultados/ablaciones_resumen.csv: knn, A11, delta_desvio
\newcommand{\deltaAucImputarModaRf}{+0{,}0011}              % resultados/ablaciones_resumen.csv: rf, A11, delta_media (ΔAUC pareado contra A0)
\newcommand{\deltaAucImputarModaRfDesvio}{0{,}0034}         % resultados/ablaciones_resumen.csv: rf, A11, delta_desvio
\newcommand{\deltaAucSinEscalarKnn}{+0{,}0084}              % resultados/ablaciones_resumen.csv: knn, A12, delta_media (ΔAUC pareado contra A0)
\newcommand{\deltaAucSinEscalarKnnDesvio}{0{,}0067}         % resultados/ablaciones_resumen.csv: knn, A12, delta_desvio
\newcommand{\pctDefaultUnknown}{20{,}8\,\%}                 % train (cargar_train()): % de filas con default = unknown
\newcommand{\pctYesDefaultUnknown}{5{,}3\,\%}               % train (cargar_train()): % de «yes» con default = unknown
\newcommand{\pctYesDefaultConocido}{12{,}8\,\%}             % train (cargar_train()): % de «yes» con default conocido
\newcommand{\filasDefaultYes}{2}                            % train (cargar_train()): filas con default = yes (D-09)
\newcommand{\mediaCampaignYes}{2{,}053}                     % train (cargar_train()): media de campaign en los «yes» (D-07)
\newcommand{\mediaCampaignNo}{2{,}627}                      % train (cargar_train()): media de campaign en los «no» (D-07)
\newcommand{\medianaCampaignYes}{2}                         % train (cargar_train()): mediana de campaign en los «yes» (D-07)
\newcommand{\medianaCampaignNo}{2}                          % train (cargar_train()): mediana de campaign en los «no» (D-07)
\newcommand{\maximoCampaign}{56}                            % train (cargar_train()): máximo de campaign (D-13)
\newcommand{\asimetriaCampaign}{4{,}89}                     % train (cargar_train()): asimetría de campaign, pandas skew (D-13)
\newcommand{\factorIqr}{1{,}5}                              % src/numeros.py: FACTOR_IQR, la regla de Tukey de resultados/eda/reporte.txt, §5 (D-17)
\newcommand{\filasAtipicasIqr}{6\,731}                      % train (cargar_train()): filas con algún valor fuera de 1,5·IQR en las 9 numéricas del modelo (D-17)
\newcommand{\pctAtipicasIqr}{20{,}4\,\%}                    % train (cargar_train()): ídem, en % de las filas (D-17)
\newcommand{\atipicosPrevious}{4\,509}                      % train (cargar_train()): filas con previous fuera de 1,5·IQR (D-17)
\newcommand{\pctYesAtipicosPrevious}{26{,}6\,\%}            % train (cargar_train()): % de «yes» entre los atípicos de previous (D-17)
\newcommand{\atipicosCampaign}{1\,911}                      % train (cargar_train()): filas con campaign fuera de 1,5·IQR (D-17)
\newcommand{\pctYesAtipicosCampaign}{4{,}5\,\%}             % train (cargar_train()): % de «yes» entre los atípicos de campaign (D-17)
\newcommand{\atipicosPdays}{1\,207}                         % train (cargar_train()): filas con pdays fuera de 1,5·IQR (D-17)
\newcommand{\pctYesAtipicosPdays}{64{,}6\,\%}               % train (cargar_train()): % de «yes» entre los atípicos de pdays (D-17)
\newcommand{\atipicosAge}{383}                              % train (cargar_train()): filas con age fuera de 1,5·IQR (D-17)
\newcommand{\pctYesAtipicosAge}{46{,}5\,\%}                 % train (cargar_train()): % de «yes» entre los atípicos de age (D-17)
\newcommand{\atipicosConsConfIdx}{357}                      % train (cargar_train()): filas con cons.conf.idx fuera de 1,5·IQR (D-17)
\newcommand{\pctYesAtipicosConsConfIdx}{40{,}6\,\%}         % train (cargar_train()): % de «yes» entre los atípicos de cons.conf.idx (D-17)

% --- Slide 8 · Métricas y presupuesto de llamadas --------------------------------------------------
\newcommand{\aucSinModelo}{0{,}500}                         % resultados/cv_final_resumen.csv: sin_modelo, validacion, auc, media
\newcommand{\aucAzar}{0{,}5}                                % resultados/cv_final_resumen.csv: sin_modelo, validacion, auc, media, con un decimal
\newcommand{\recallSinModelo}{0{,}200}                      % resultados/cv_final_resumen.csv: sin_modelo, validacion, recall_q, media
\newcommand{\precisionSinModelo}{0{,}113}                   % resultados/cv_final_resumen.csv: sin_modelo, validacion, precision_q, media
\newcommand{\apSinModelo}{0{,}113}                          % resultados/cv_final_resumen.csv: sin_modelo, validacion, ap, media
\newcommand{\fUnoSinModelo}{0{,}144}                        % resultados/cv_final_resumen.csv: sin_modelo, validacion, f1_q, media
\newcommand{\presupuesto}{20\,\%}                           % src/metricas.py: PRESUPUESTO (D-20)
\newcommand{\presupuestoBajo}{10\,\%}                       % resultados/sensibilidad_q_final.csv: el menor q
\newcommand{\presupuestoAlto}{30\,\%}                       % resultados/sensibilidad_q_final.csv: el mayor q
\newcommand{\rotuloLineaBase}{sin modelo}                   % src/estilo.py: ROTULO_LINEA_BASE (texto)

% --- Slide 9 · Los modelos sin ajustar -------------------------------------------------------------
\newcommand{\aucNbReferencia}{0{,}782}                      % resultados/cv_referencia_resumen.csv: nb_categorico, validacion, auc, media
\newcommand{\aucNbReferenciaDesvio}{0{,}007}                % resultados/cv_referencia_resumen.csv: nb_categorico, validacion, auc, desvio
\newcommand{\recallNbReferencia}{0{,}614}                   % resultados/cv_referencia_resumen.csv: nb_categorico, validacion, recall_q, media
\newcommand{\aucTrainNbReferencia}{0{,}783}                 % resultados/cv_referencia_resumen.csv: nb_categorico, train, auc, media
\newcommand{\brechaNbReferencia}{0{,}001}                   % resultados/cv_referencia_resumen.csv: nb_categorico, brecha, auc, media (train − validación)
\newcommand{\aucNbGaussianoVal}{0{,}769}                    % resultados/cv_referencia_resumen.csv: nb_gaussiano, validacion, auc, media
\newcommand{\aucNbGaussianoValDesvio}{0{,}009}              % resultados/cv_referencia_resumen.csv: nb_gaussiano, validacion, auc, desvio
\newcommand{\recallNbGaussianoVal}{0{,}566}                 % resultados/cv_referencia_resumen.csv: nb_gaussiano, validacion, recall_q, media
\newcommand{\aucTrainNbGaussianoVal}{0{,}771}               % resultados/cv_referencia_resumen.csv: nb_gaussiano, train, auc, media
\newcommand{\brechaNbGaussianoVal}{0{,}002}                 % resultados/cv_referencia_resumen.csv: nb_gaussiano, brecha, auc, media (train − validación)
\newcommand{\aucSvmRbfReferencia}{0{,}702}                  % resultados/cv_referencia_resumen.csv: svm, validacion, auc, media
\newcommand{\aucSvmRbfReferenciaDesvio}{0{,}009}            % resultados/cv_referencia_resumen.csv: svm, validacion, auc, desvio
\newcommand{\recallSvmRbfReferencia}{0{,}549}               % resultados/cv_referencia_resumen.csv: svm, validacion, recall_q, media
\newcommand{\aucTrainSvmRbfReferencia}{0{,}901}             % resultados/cv_referencia_resumen.csv: svm, train, auc, media
\newcommand{\brechaSvmRbfReferencia}{0{,}200}               % resultados/cv_referencia_resumen.csv: svm, brecha, auc, media (train − validación)
\newcommand{\aucKnnReferencia}{0{,}755}                     % resultados/cv_referencia_resumen.csv: knn, validacion, auc, media
\newcommand{\aucKnnReferenciaDesvio}{0{,}012}               % resultados/cv_referencia_resumen.csv: knn, validacion, auc, desvio
\newcommand{\recallKnnReferencia}{0{,}593}                  % resultados/cv_referencia_resumen.csv: knn, validacion, recall_q, media
\newcommand{\aucTrainKnnReferencia}{0{,}871}                % resultados/cv_referencia_resumen.csv: knn, train, auc, media
\newcommand{\brechaKnnReferencia}{0{,}116}                  % resultados/cv_referencia_resumen.csv: knn, brecha, auc, media (train − validación)
\newcommand{\aucRfReferencia}{0{,}772}                      % resultados/cv_referencia_resumen.csv: rf, validacion, auc, media
\newcommand{\aucRfReferenciaDesvio}{0{,}006}                % resultados/cv_referencia_resumen.csv: rf, validacion, auc, desvio
\newcommand{\recallRfReferencia}{0{,}605}                   % resultados/cv_referencia_resumen.csv: rf, validacion, recall_q, media
\newcommand{\aucTrainRfReferencia}{1{,}000}                 % resultados/cv_referencia_resumen.csv: rf, train, auc, media
\newcommand{\brechaRfReferencia}{0{,}228}                   % resultados/cv_referencia_resumen.csv: rf, brecha, auc, media (train − validación)
\newcommand{\modeloMejorReferencia}{Naive Bayes}            % resultados/cv_referencia_resumen.csv: el de mayor AUC de validación entre los cuatro del enunciado (texto)
\newcommand{\margenAucReferencia}{0{,}282}                  % resultados/cv_referencia_resumen.csv: el mayor AUC de validación de los cuatro menos el de sin_modelo
\newcommand{\deltaAucNbCategorico}{+0{,}013}                % resultados/cv_referencia.csv: nb_categorico − nb_gaussiano, AUC de validación fold a fold, media (D-18)
\newcommand{\deltaAucNbCategoricoDesvio}{0{,}0024}          % resultados/cv_referencia.csv: ídem, desvío
\newcommand{\deltaAucNbCategoricoMinimo}{+0{,}010}          % resultados/cv_referencia.csv: ídem, el menor fold
\newcommand{\deltaAucNbCategoricoMaximo}{+0{,}016}          % resultados/cv_referencia.csv: ídem, el mayor fold
\newcommand{\deltaRecallNbCategorico}{+0{,}048}             % resultados/cv_referencia.csv: nb_categorico − nb_gaussiano, recall_q de validación fold a fold, media (D-18)
\newcommand{\deltaRecallNbCategoricoDesvio}{0{,}0083}       % resultados/cv_referencia.csv: ídem, desvío
\newcommand{\vecinosKnnReferencia}{15}                      % resultados/cv_referencia_resumen.csv: knn, configuracion, n_neighbors
\newcommand{\arbolesRfReferencia}{300}                      % resultados/cv_referencia_resumen.csv: rf, configuracion, n_estimators
\newcommand{\cSvmReferencia}{1}                             % resultados/cv_referencia_resumen.csv: svm, configuracion, C
\newcommand{\cortesNb}{10}                                  % resultados/cv_final_resumen.csv: nb_categorico, configuracion, n_cortes (D-18, D-22)
\newcommand{\alfaNb}{1}                                     % resultados/cv_final_resumen.csv: nb_categorico, configuracion, alpha, la corrección de Laplace (D-18)

% --- Slide 10 · Curva de RF ------------------------------------------------------------------------
\newcommand{\profundidadRf}{8}                              % resultados/cv_final_resumen.csv: rf, configuracion, max_depth
\newcommand{\arbolesRf}{200}                                % resultados/cv_final_resumen.csv: rf, configuracion, n_estimators
\newcommand{\aucRfProfundidadElegida}{0{,}795}              % resultados/hiperparametros.json: curvas.rf_max_depth.auc_validacion_media
\newcommand{\errorEstandarRfProfundidadElegida}{0{,}0025}   % resultados/hiperparametros.json: curvas.rf_max_depth.error_estandar
\newcommand{\brechaRfProfundidadElegida}{0{,}030}           % resultados/hiperparametros.json: curvas.rf_max_depth.brecha
\newcommand{\profundidadRfMejor}{10}                        % resultados/hiperparametros.json: curvas.rf_max_depth.mejor
\newcommand{\aucRfProfundidadMejor}{0{,}797}                % resultados/hiperparametros.json: curvas.rf_max_depth.auc_validacion_mejor
\newcommand{\umbralUnoEsRfProfundidad}{0{,}7939}            % resultados/hiperparametros.json: curvas.rf_max_depth.umbral_1es
\newcommand{\errorEstandarRfProfundidadMejor}{0{,}0029}     % resultados/hiperparametros.json: curvas.rf_max_depth.error_estandar_mejor
\newcommand{\profundidadRfMinima}{2}                        % src/modelos.py: GRILLAS, la menor max_depth de la grilla de RF (el extremo del subajuste)
\newcommand{\arbolesRfRegla}{25}                            % resultados/hiperparametros.json: curvas.rf_n_estimators_depth8.valor, lo que elegía la regla de 1 ES (N0-11)
\newcommand{\arbolesRfMejor}{800}                           % resultados/hiperparametros.json: curvas.rf_n_estimators_depth8.mejor (N0-11)
\newcommand{\umbralUnoEsRfArboles}{0{,}7930}                % resultados/hiperparametros.json: curvas.rf_n_estimators_depth8.umbral_1es (N0-11)
\newcommand{\aucRfProfundidadDos}{0{,}782}                  % resultados/curvas/rf_max_depth.csv: punto 2, validacion, auc, media
\newcommand{\aucRfProfundidadCuatro}{0{,}789}               % resultados/curvas/rf_max_depth.csv: punto 4, validacion, auc, media
\newcommand{\aucRfProfundidadSeis}{0{,}794}                 % resultados/curvas/rf_max_depth.csv: punto 6, validacion, auc, media
\newcommand{\aucRfProfundidadOcho}{0{,}795}                 % resultados/curvas/rf_max_depth.csv: punto 8, validacion, auc, media
\newcommand{\aucRfProfundidadDiez}{0{,}797}                 % resultados/curvas/rf_max_depth.csv: punto 10, validacion, auc, media
\newcommand{\aucRfProfundidadDoce}{0{,}795}                 % resultados/curvas/rf_max_depth.csv: punto 12, validacion, auc, media
\newcommand{\aucRfProfundidadQuince}{0{,}793}               % resultados/curvas/rf_max_depth.csv: punto 15, validacion, auc, media
\newcommand{\aucRfProfundidadVeinte}{0{,}784}               % resultados/curvas/rf_max_depth.csv: punto 20, validacion, auc, media
\newcommand{\aucRfProfundidadVeinticinco}{0{,}776}          % resultados/curvas/rf_max_depth.csv: punto 25, validacion, auc, media
\newcommand{\aucRfProfundidadSinLimite}{0{,}772}            % resultados/curvas/rf_max_depth.csv: punto None, validacion, auc, media
\newcommand{\aucRfProfundidadDosConCuatroDecimales}{0{,}7821} % resultados/curvas/rf_max_depth.csv: punto 2, validacion, auc, media, con cuatro decimales (para compararla con el umbral de 1 ES)
\newcommand{\aucRfProfundidadCuatroConCuatroDecimales}{0{,}7893} % resultados/curvas/rf_max_depth.csv: punto 4, validacion, auc, media, con cuatro decimales (para compararla con el umbral de 1 ES)
\newcommand{\aucRfProfundidadSeisConCuatroDecimales}{0{,}7936} % resultados/curvas/rf_max_depth.csv: punto 6, validacion, auc, media, con cuatro decimales (para compararla con el umbral de 1 ES)
\newcommand{\aucRfProfundidadOchoConCuatroDecimales}{0{,}7954} % resultados/curvas/rf_max_depth.csv: punto 8, validacion, auc, media, con cuatro decimales (para compararla con el umbral de 1 ES)
\newcommand{\aucRfProfundidadDiezConCuatroDecimales}{0{,}7968} % resultados/curvas/rf_max_depth.csv: punto 10, validacion, auc, media, con cuatro decimales (para compararla con el umbral de 1 ES)
\newcommand{\aucRfProfundidadDoceConCuatroDecimales}{0{,}7952} % resultados/curvas/rf_max_depth.csv: punto 12, validacion, auc, media, con cuatro decimales (para compararla con el umbral de 1 ES)
\newcommand{\aucRfProfundidadQuinceConCuatroDecimales}{0{,}7931} % resultados/curvas/rf_max_depth.csv: punto 15, validacion, auc, media, con cuatro decimales (para compararla con el umbral de 1 ES)
\newcommand{\aucRfProfundidadVeinteConCuatroDecimales}{0{,}7838} % resultados/curvas/rf_max_depth.csv: punto 20, validacion, auc, media, con cuatro decimales (para compararla con el umbral de 1 ES)
\newcommand{\aucRfProfundidadVeinticincoConCuatroDecimales}{0{,}7757} % resultados/curvas/rf_max_depth.csv: punto 25, validacion, auc, media, con cuatro decimales (para compararla con el umbral de 1 ES)
\newcommand{\aucRfProfundidadSinLimiteConCuatroDecimales}{0{,}7720} % resultados/curvas/rf_max_depth.csv: punto None, validacion, auc, media, con cuatro decimales (para compararla con el umbral de 1 ES)
\newcommand{\aucTrainRfProfundidadSinLimite}{1{,}000}       % resultados/curvas/rf_max_depth.csv: punto None, train, auc, media
\newcommand{\brechaRfProfundidadSinLimite}{0{,}228}         % resultados/curvas/rf_max_depth.csv: punto None, train − validación
\newcommand{\aucTrainRfProfundidadDos}{0{,}783}             % resultados/curvas/rf_max_depth.csv: punto 2, train, auc, media
\newcommand{\brechaRfProfundidadDos}{0{,}001}               % resultados/curvas/rf_max_depth.csv: punto 2, train − validación
\newcommand{\aucRfArbolesDiez}{0{,}791}                     % resultados/curvas/rf_n_estimators_depth8.csv: punto 10, validacion, auc, media (N0-11)
\newcommand{\aucRfArbolesVeinticinco}{0{,}793}              % resultados/curvas/rf_n_estimators_depth8.csv: punto 25, validacion, auc, media (N0-11)
\newcommand{\aucRfArbolesCincuenta}{0{,}795}                % resultados/curvas/rf_n_estimators_depth8.csv: punto 50, validacion, auc, media (N0-11)
\newcommand{\aucRfArbolesCien}{0{,}795}                     % resultados/curvas/rf_n_estimators_depth8.csv: punto 100, validacion, auc, media (N0-11)
\newcommand{\aucRfArbolesDoscientos}{0{,}795}               % resultados/curvas/rf_n_estimators_depth8.csv: punto 200, validacion, auc, media (N0-11)
\newcommand{\aucRfArbolesCuatrocientos}{0{,}796}            % resultados/curvas/rf_n_estimators_depth8.csv: punto 400, validacion, auc, media (N0-11)
\newcommand{\aucRfArbolesOchocientos}{0{,}796}              % resultados/curvas/rf_n_estimators_depth8.csv: punto 800, validacion, auc, media (N0-11)
\newcommand{\aucRfSinPesos}{0{,}795}                        % resultados/curvas/rf_pesos_clase_depth8.csv: punto None, validacion, auc, media (N0-12)
\newcommand{\aucRfPesosBalanceados}{0{,}796}                % resultados/curvas/rf_pesos_clase_depth8.csv: punto balanced, validacion, auc, media (N0-12)
\newcommand{\deltaAucPesosRf}{+0{,}0006}                    % resultados/curvas/rf_pesos_clase_depth8.csv: balanced − None (N0-12)
\newcommand{\errorEstandarRfPesosBalanceados}{0{,}0024}     % resultados/curvas/rf_pesos_clase_depth8.csv: punto balanced, error estándar (N0-12)
\newcommand{\deltaAucPesosRfDesvio}{0{,}0009}               % resultados/curvas/rf_pesos_clase_depth8.csv: balanced − None fold a fold, desvío; la media es \deltaAucPesosRf (N0-12)
\newcommand{\foldsPositivosPesosRf}{5}                      % resultados/curvas/rf_pesos_clase_depth8.csv: folds en que balanced − None es positiva (N0-12)

% --- Slide 11 · Curva de KNN -----------------------------------------------------------------------
\newcommand{\vecinosKnn}{801}                               % resultados/cv_final_resumen.csv: knn, configuracion, n_neighbors
\newcommand{\vecinosKnnMejor}{601}                          % resultados/hiperparametros.json: curvas.knn_n_neighbors_uniform.mejor
\newcommand{\aucKnnMejor}{0{,}784}                          % resultados/hiperparametros.json: curvas.knn_n_neighbors_uniform.auc_validacion_mejor
\newcommand{\vecinosKnnMinimoUnoEs}{151}                    % resultados/hiperparametros.json: el menor de curvas.knn_n_neighbors_uniform.puntos_dentro_1es
\newcommand{\umbralUnoEsKnn}{0{,}7806}                      % resultados/hiperparametros.json: curvas.knn_n_neighbors_uniform.umbral_1es
\newcommand{\errorEstandarKnnMejor}{0{,}0037}               % resultados/hiperparametros.json: curvas.knn_n_neighbors_uniform.error_estandar_mejor
\newcommand{\vecinosKnnMinimo}{1}                           % src/modelos.py: GRILLAS, el menor n_neighbors de la grilla de KNN (el extremo del sobreajuste)
\newcommand{\vecinosKnnMeseta}{301}                         % DECISIONES.md, D-22: el k desde el que la curva de KNN es una meseta (una lectura de la curva, no una regla); «?» si ya no está dentro de 1 ES
\newcommand{\aucKnnUno}{0{,}619}                            % resultados/curvas/knn_n_neighbors_uniform.csv: punto 1, validacion, auc, media
\newcommand{\aucKnnTres}{0{,}696}                           % resultados/curvas/knn_n_neighbors_uniform.csv: punto 3, validacion, auc, media
\newcommand{\aucKnnCinco}{0{,}724}                          % resultados/curvas/knn_n_neighbors_uniform.csv: punto 5, validacion, auc, media
\newcommand{\aucKnnNueve}{0{,}741}                          % resultados/curvas/knn_n_neighbors_uniform.csv: punto 9, validacion, auc, media
\newcommand{\aucKnnQuince}{0{,}755}                         % resultados/curvas/knn_n_neighbors_uniform.csv: punto 15, validacion, auc, media
\newcommand{\aucKnnVeinticinco}{0{,}765}                    % resultados/curvas/knn_n_neighbors_uniform.csv: punto 25, validacion, auc, media
\newcommand{\aucKnnCuarentaYUno}{0{,}769}                   % resultados/curvas/knn_n_neighbors_uniform.csv: punto 41, validacion, auc, media
\newcommand{\aucKnnSesentaYUno}{0{,}773}                    % resultados/curvas/knn_n_neighbors_uniform.csv: punto 61, validacion, auc, media
\newcommand{\aucKnnCientoUno}{0{,}779}                      % resultados/curvas/knn_n_neighbors_uniform.csv: punto 101, validacion, auc, media
\newcommand{\aucKnnCientoCincuentaYUno}{0{,}782}            % resultados/curvas/knn_n_neighbors_uniform.csv: punto 151, validacion, auc, media
\newcommand{\aucKnnDoscientosUno}{0{,}783}                  % resultados/curvas/knn_n_neighbors_uniform.csv: punto 201, validacion, auc, media
\newcommand{\aucKnnTrescientosUno}{0{,}783}                 % resultados/curvas/knn_n_neighbors_uniform.csv: punto 301, validacion, auc, media
\newcommand{\aucKnnCuatrocientosUno}{0{,}784}               % resultados/curvas/knn_n_neighbors_uniform.csv: punto 401, validacion, auc, media
\newcommand{\aucKnnSeiscientosUno}{0{,}784}                 % resultados/curvas/knn_n_neighbors_uniform.csv: punto 601, validacion, auc, media
\newcommand{\aucKnnOchocientosUno}{0{,}784}                 % resultados/curvas/knn_n_neighbors_uniform.csv: punto 801, validacion, auc, media
\newcommand{\aucTrainKnnUno}{0{,}988}                       % resultados/curvas/knn_n_neighbors_uniform.csv: punto 1, train, auc, media
\newcommand{\brechaKnnUno}{0{,}369}                         % resultados/curvas/knn_n_neighbors_uniform.csv: punto 1, train − validación
\newcommand{\vecinosKnnDistancia}{801}                      % resultados/curvas/knn_n_neighbors_distance.csv: el punto que elige D-22 sobre la curva (src/curvas.py, elegir_punto; weights = distance), n_neighbors
\newcommand{\aucKnnDistancia}{0{,}770}                      % resultados/curvas/knn_n_neighbors_distance.csv: el punto que elige D-22 sobre la curva (src/curvas.py, elegir_punto; weights = distance), validacion, auc, media
\newcommand{\aucTrainKnnDistancia}{1{,}000}                 % resultados/curvas/knn_n_neighbors_distance.csv: el punto que elige D-22 sobre la curva (src/curvas.py, elegir_punto; weights = distance), train, auc, media
\newcommand{\brechaKnnDistancia}{0{,}230}                   % resultados/curvas/knn_n_neighbors_distance.csv: el punto que elige D-22 sobre la curva (src/curvas.py, elegir_punto; weights = distance), train − validación

% --- Slide 12 · Curva de la SVM --------------------------------------------------------------------
\newcommand{\cSvm}{0{,}001}                                 % resultados/cv_final_resumen.csv: svm, configuracion, C
\newcommand{\kernelSvm}{lineal}                             % resultados/cv_final_resumen.csv: svm, configuracion, kernel (texto)
\newcommand{\aucSvmBalanceada}{0{,}770}                     % resultados/curvas/svm_pesos_clase.csv: punto balanced, validacion, auc, media (RBF, C = 1)
\newcommand{\deltaAucPesosSvm}{+0{,}068}                    % resultados/curvas/svm_pesos_clase.csv: balanced − None (D-21, N0-12)
\newcommand{\deltaAucPesosSvmDesvio}{0{,}011}               % resultados/curvas/svm_pesos_clase.csv: balanced − None fold a fold, desvío; la media es \deltaAucPesosSvm (D-21)
\newcommand{\foldsPositivosPesosSvm}{5}                     % resultados/curvas/svm_pesos_clase.csv: folds en que balanced − None es positiva (D-21)
\newcommand{\cSvmSinPesosMejor}{0{,}1}                      % resultados/hiperparametros.json: curvas.svm_C.mejor (RBF, sin pesos)
\newcommand{\aucSvmSinPesosMejor}{0{,}706}                  % resultados/hiperparametros.json: curvas.svm_C.auc_validacion_mejor
\newcommand{\cSvmBalanceadaMejor}{0{,}01}                   % resultados/hiperparametros.json: curvas.svm_C_balanced.mejor (RBF, pesos balanceados)
\newcommand{\aucSvmBalanceadaMejor}{0{,}772}                % resultados/hiperparametros.json: curvas.svm_C_balanced.auc_validacion_mejor
\newcommand{\cSvmKernel}{0{,}001}                           % resultados/curvas/svm_kernel.csv: C de la configuración en que se comparan los kernels
\newcommand{\aucSvmKernelLineal}{0{,}774}                   % resultados/curvas/svm_kernel.csv: punto linear, validacion, auc, media
\newcommand{\aucSvmKernelPolinomico}{0{,}774}               % resultados/curvas/svm_kernel.csv: punto poly, validacion, auc, media
\newcommand{\aucSvmKernelRbf}{0{,}768}                      % resultados/curvas/svm_kernel.csv: punto rbf, validacion, auc, media
\newcommand{\cSvmLinealMejor}{0{,}01}                       % resultados/hiperparametros.json: curvas.svm_C_linear_balanced.mejor
\newcommand{\aucSvmLinealMejor}{0{,}774}                    % resultados/hiperparametros.json: curvas.svm_C_linear_balanced.auc_validacion_mejor
\newcommand{\segundosLimiteCosto}{600}                      % resultados/costos.csv: el límite de la corrida excedida (s)
\newcommand{\cSvmCostoExcedido}{10}                         % resultados/costos.csv: C de la corrida excedida (SVC lineal)

% --- Slide 13 · El modelo final en validación ------------------------------------------------------
\newcommand{\aucRfVal}{0{,}795}                             % resultados/cv_final_resumen.csv: rf, validacion, auc, media
\newcommand{\aucRfValDesvio}{0{,}006}                       % resultados/cv_final_resumen.csv: rf, validacion, auc, desvio
\newcommand{\recallRfVal}{0{,}629}                          % resultados/cv_final_resumen.csv: rf, validacion, recall_q, media
\newcommand{\recallRfValDesvio}{0{,}018}                    % resultados/cv_final_resumen.csv: rf, validacion, recall_q, desvio
\newcommand{\precisionRfVal}{0{,}354}                       % resultados/cv_final_resumen.csv: rf, validacion, precision_q, media
\newcommand{\apRfVal}{0{,}462}                              % resultados/cv_final_resumen.csv: rf, validacion, ap, media
\newcommand{\fUnoRfVal}{0{,}454}                            % resultados/cv_final_resumen.csv: rf, validacion, f1_q, media
\newcommand{\aucTrainRf}{0{,}825}                           % resultados/cv_final_resumen.csv: rf, train, auc, media
\newcommand{\brechaRf}{0{,}030}                             % resultados/cv_final_resumen.csv: rf, brecha, auc, media (train − validación)
\newcommand{\aucKnnVal}{0{,}784}                            % resultados/cv_final_resumen.csv: knn, validacion, auc, media
\newcommand{\aucKnnValDesvio}{0{,}008}                      % resultados/cv_final_resumen.csv: knn, validacion, auc, desvio
\newcommand{\recallKnnVal}{0{,}617}                         % resultados/cv_final_resumen.csv: knn, validacion, recall_q, media
\newcommand{\recallKnnValDesvio}{0{,}016}                   % resultados/cv_final_resumen.csv: knn, validacion, recall_q, desvio
\newcommand{\precisionKnnVal}{0{,}348}                      % resultados/cv_final_resumen.csv: knn, validacion, precision_q, media
\newcommand{\apKnnVal}{0{,}429}                             % resultados/cv_final_resumen.csv: knn, validacion, ap, media
\newcommand{\fUnoKnnVal}{0{,}445}                           % resultados/cv_final_resumen.csv: knn, validacion, f1_q, media
\newcommand{\aucTrainKnn}{0{,}791}                          % resultados/cv_final_resumen.csv: knn, train, auc, media
\newcommand{\brechaKnn}{0{,}006}                            % resultados/cv_final_resumen.csv: knn, brecha, auc, media (train − validación)
\newcommand{\aucNbVal}{0{,}782}                             % resultados/cv_final_resumen.csv: nb_categorico, validacion, auc, media
\newcommand{\aucNbValDesvio}{0{,}007}                       % resultados/cv_final_resumen.csv: nb_categorico, validacion, auc, desvio
\newcommand{\recallNbVal}{0{,}614}                          % resultados/cv_final_resumen.csv: nb_categorico, validacion, recall_q, media
\newcommand{\recallNbValDesvio}{0{,}017}                    % resultados/cv_final_resumen.csv: nb_categorico, validacion, recall_q, desvio
\newcommand{\precisionNbVal}{0{,}346}                       % resultados/cv_final_resumen.csv: nb_categorico, validacion, precision_q, media
\newcommand{\apNbVal}{0{,}422}                              % resultados/cv_final_resumen.csv: nb_categorico, validacion, ap, media
\newcommand{\fUnoNbVal}{0{,}443}                            % resultados/cv_final_resumen.csv: nb_categorico, validacion, f1_q, media
\newcommand{\aucTrainNb}{0{,}783}                           % resultados/cv_final_resumen.csv: nb_categorico, train, auc, media
\newcommand{\brechaNb}{0{,}001}                             % resultados/cv_final_resumen.csv: nb_categorico, brecha, auc, media (train − validación)
\newcommand{\aucSvmVal}{0{,}774}                            % resultados/cv_final_resumen.csv: svm, validacion, auc, media
\newcommand{\aucSvmValDesvio}{0{,}009}                      % resultados/cv_final_resumen.csv: svm, validacion, auc, desvio
\newcommand{\recallSvmVal}{0{,}614}                         % resultados/cv_final_resumen.csv: svm, validacion, recall_q, media
\newcommand{\recallSvmValDesvio}{0{,}018}                   % resultados/cv_final_resumen.csv: svm, validacion, recall_q, desvio
\newcommand{\precisionSvmVal}{0{,}346}                      % resultados/cv_final_resumen.csv: svm, validacion, precision_q, media
\newcommand{\apSvmVal}{0{,}405}                             % resultados/cv_final_resumen.csv: svm, validacion, ap, media
\newcommand{\fUnoSvmVal}{0{,}442}                           % resultados/cv_final_resumen.csv: svm, validacion, f1_q, media
\newcommand{\aucTrainSvm}{0{,}776}                          % resultados/cv_final_resumen.csv: svm, train, auc, media
\newcommand{\brechaSvm}{0{,}002}                            % resultados/cv_final_resumen.csv: svm, brecha, auc, media (train − validación)
\newcommand{\deltaAucAjusteRf}{+0{,}023}                    % resultados/cv_final_resumen.csv menos resultados/cv_referencia_resumen.csv: rf, validacion, auc, media
\newcommand{\deltaAucAjusteKnn}{+0{,}029}                   % resultados/cv_final_resumen.csv menos resultados/cv_referencia_resumen.csv: knn, validacion, auc, media
\newcommand{\deltaAucAjusteSvm}{+0{,}072}                   % resultados/cv_final_resumen.csv menos resultados/cv_referencia_resumen.csv: svm, validacion, auc, media
\newcommand{\modeloFinal}{Random Forest}                    % resultados/modelo_elegido.json: ranking, 1.º (texto)
\newcommand{\modeloSegundoFinal}{KNN}                       % resultados/modelo_elegido.json: ranking, 2.º (texto)
\newcommand{\modeloUltimoFinal}{SVM}                        % resultados/modelo_elegido.json: ranking, último (texto)
\newcommand{\errorEstandarAucRf}{0{,}0025}                  % resultados/cv_final_resumen.csv: rf, validacion, auc, error_estandar
\newcommand{\umbralEmpateFinal}{0{,}7927}                   % resultados/modelo_elegido.json: ranking, AUC del 1.º menos su error estándar (D-23)
\newcommand{\modelosDentroUnoEs}{0}                         % resultados/modelo_elegido.json: cuántos quedan dentro de 1 ES del mejor (empate_dentro_1es)
\newcommand{\margenAucFinal}{0{,}011}                       % resultados/modelo_elegido.json: ranking, AUC del 1.º menos el del 2.º
\newcommand{\margenAucSinModeloFinal}{0{,}295}              % resultados/modelo_elegido.json: ranking, AUC del 1.º menos linea_base.auc
\newcommand{\rangoAucFinal}{0{,}021}                        % resultados/modelo_elegido.json: ranking, AUC del 1.º menos el del último

% --- Slides 14 y 15 · Test y matriz de confusión ---------------------------------------------------
\newcommand{\yesTest}{?}                                    % resultados/evaluacion_test.json: n_yes_test
\newcommand{\pctYesTest}{?}                                 % resultados/evaluacion_test.json: n_yes_test / n_test
\newcommand{\diferenciaAucTestValidacion}{?}                % resultados/evaluacion_test.json: metricas.auc.valor menos resultados/modelo_elegido.json: metricas.validacion.auc.media (sólo en voz alta, guía C2)
\newcommand{\diferenciaRecallTestValidacion}{?}             % resultados/evaluacion_test.json: metricas.recall_q.valor menos resultados/modelo_elegido.json: metricas.validacion.recall_q.media (sólo en voz alta, guía C2)
\newcommand{\vpRfVal}{2\,329}                               % resultados/oof_final_rf.csv con la y de train (cargar_train()): VP llamando al 20 % de la lista fuera de fold
\newcommand{\fpRfVal}{4\,259}                               % resultados/oof_final_rf.csv con la y de train (cargar_train()): FP llamando al 20 % de la lista fuera de fold
\newcommand{\fnRfVal}{1\,382}                               % resultados/oof_final_rf.csv con la y de train (cargar_train()): FN llamando al 20 % de la lista fuera de fold
\newcommand{\vnRfVal}{24\,970}                              % resultados/oof_final_rf.csv con la y de train (cargar_train()): VN llamando al 20 % de la lista fuera de fold
\newcommand{\llamadasPorYesRf}{2{,}8}                       % resultados/oof_final_rf.csv con la y de train (cargar_train()): llamadas por cada VP, llamando al 20 % de la lista fuera de fold, (VP + FP) / VP
\newcommand{\llamadasPorYesSinModelo}{8{,}9}                % resultados/cv_final_resumen.csv: 1 / (sin_modelo, validacion, precision_q, media), las llamadas por cada «yes» sin modelo
\newcommand{\nivelIntervalo}{95\,\%}                        % src/evaluar_test.py: PERCENTILES, el nivel del intervalo por bootstrap del AUC y del recall de test (D-24)
\newcommand{\remuestreosBootstrap}{2\,000}                  % src/evaluar_test.py: N_BOOTSTRAP, los remuestreos de ese intervalo (D-24)

% --- Slide 16 · Por qué cada modelo rindió lo que rindió -------------------------------------------
\newcommand{\correlacionMacroMinima}{0{,}91}                % train (cargar_train()): el menor |r| entre los pares del bloque macro con |r| >= 0,9 (D-14)
\newcommand{\correlacionMacroMaxima}{0{,}97}                % train (cargar_train()): el mayor |r| entre esos pares (D-14)
\newcommand{\desvioPdays}{186{,}6}                          % train (cargar_train()): desvío estándar de pdays (D-12)
\newcommand{\desvioNrEmployed}{72{,}4}                      % train (cargar_train()): desvío estándar de nr.employed (D-12)
\newcommand{\pctVarianzaPdays}{86{,}6\,\%}                  % train (cargar_train()): % de la varianza de las 9 numéricas del modelo que es de pdays (D-12)
\newcommand{\pctVarianzaNrEmployed}{13{,}0\,\%}             % train (cargar_train()): ídem, de nr.employed (D-12)
\newcommand{\razonDesviosNumericas}{374}                    % train (cargar_train()): el mayor desvío de las 9 numéricas del modelo / el menor
\newcommand{\probabilidadAltaNb}{0{,}99}                    % src/numeros.py: PROBABILIDAD_ALTA_NB, la P(yes) de Naive Bayes que se cuenta como «casi 1» (una elección)
\newcommand{\filasNbProbabilidadAlta}{3\,121}               % resultados/oof_final_nb_categorico.csv con la y de train (cargar_train()): filas con P(yes) >= PROBABILIDAD_ALTA_NB fuera de fold, P(yes) = 1 / (1 + e^(−puntaje)), porque el puntaje es el log-odds
\newcommand{\pctNbProbabilidadAlta}{9{,}5\,\%}              % resultados/oof_final_nb_categorico.csv con la y de train (cargar_train()): ídem, en % de las filas de train
\newcommand{\pctYesNbProbabilidadAlta}{47{,}3\,\%}          % resultados/oof_final_nb_categorico.csv con la y de train (cargar_train()): % de «yes» entre esas filas
\newcommand{\probabilidadMediaNbDecilSuperior}{0{,}998}     % resultados/oof_final_nb_categorico.csv con la y de train (cargar_train()): P(yes) media en el 10 % de la lista con mayor puntaje (empates del corte en proporción, como src/metricas.py)
\newcommand{\tasaYesNbDecilSuperior}{0{,}464}               % resultados/oof_final_nb_categorico.csv con la y de train (cargar_train()): proporción de «yes» en ese mismo 10 %

% --- Slide 17 · Limitaciones: sensibilidad al presupuesto ------------------------------------------
\newcommand{\recallRfDiez}{0{,}446}                         % resultados/sensibilidad_q_final.csv: rf, q = 0.1, recall_q, media de los folds
\newcommand{\recallRfTreinta}{0{,}705}                      % resultados/sensibilidad_q_final.csv: rf, q = 0.3, recall_q, media de los folds
\newcommand{\recallKnnDiez}{0{,}414}                        % resultados/sensibilidad_q_final.csv: knn, q = 0.1, recall_q, media de los folds
\newcommand{\recallKnnTreinta}{0{,}690}                     % resultados/sensibilidad_q_final.csv: knn, q = 0.3, recall_q, media de los folds
\newcommand{\recallNbDiez}{0{,}411}                         % resultados/sensibilidad_q_final.csv: nb_categorico, q = 0.1, recall_q, media de los folds
\newcommand{\recallNbTreinta}{0{,}697}                      % resultados/sensibilidad_q_final.csv: nb_categorico, q = 0.3, recall_q, media de los folds
\newcommand{\recallSvmDiez}{0{,}415}                        % resultados/sensibilidad_q_final.csv: svm, q = 0.1, recall_q, media de los folds
\newcommand{\recallSvmTreinta}{0{,}695}                     % resultados/sensibilidad_q_final.csv: svm, q = 0.3, recall_q, media de los folds
\newcommand{\recallSinModeloDiez}{0{,}100}                  % resultados/sensibilidad_q_final.csv: sin_modelo, q = 0.1, recall_q, media de los folds
\newcommand{\recallSinModeloTreinta}{0{,}300}               % resultados/sensibilidad_q_final.csv: sin_modelo, q = 0.3, recall_q, media de los folds
\newcommand{\precisionRfDiez}{0{,}502}                      % resultados/sensibilidad_q_final.csv: rf, q = 0.1, precision_q, media de los folds
\newcommand{\precisionRfVeinte}{0{,}354}                    % resultados/sensibilidad_q_final.csv: rf, q = 0.2, precision_q, media de los folds
\newcommand{\precisionRfTreinta}{0{,}265}                   % resultados/sensibilidad_q_final.csv: rf, q = 0.3, precision_q, media de los folds

% --- Slides 18 y 19 · Hallazgo: el modelo aprende la época -----------------------------------------
\newcommand{\aucRfAdelante}{0{,}558}                        % resultados/robustez_temporal_resumen.csv: rf, hacia_adelante, todas, auc, media
\newcommand{\aucRfAdelanteDesvio}{0{,}123}                  % resultados/robustez_temporal_resumen.csv: rf, hacia_adelante, todas, auc, desvio
\newcommand{\aucRfSinMacroAdelante}{0{,}556}                % resultados/robustez_temporal_resumen.csv: rf, hacia_adelante, sin_macro, auc, media
\newcommand{\aucRfSinMacroAdelanteDesvio}{0{,}110}          % resultados/robustez_temporal_resumen.csv: rf, hacia_adelante, sin_macro, auc, desvio
\newcommand{\aucRfSinMacroBarajado}{0{,}772}                % resultados/robustez_temporal_resumen.csv: rf, barajado, sin_macro, auc, media
\newcommand{\aucRfSinMacroBarajadoDesvio}{0{,}010}          % resultados/robustez_temporal_resumen.csv: rf, barajado, sin_macro, auc, desvio
\newcommand{\aucRfSinMacroNiMonthAdelante}{0{,}547}         % resultados/robustez_temporal_resumen.csv: rf, hacia_adelante, sin_macro_ni_month, auc, media
\newcommand{\aucRfSinMacroNiMonthAdelanteDesvio}{0{,}098}   % resultados/robustez_temporal_resumen.csv: rf, hacia_adelante, sin_macro_ni_month, auc, desvio
\newcommand{\aucRfSinMacroNiMonthBarajado}{0{,}736}         % resultados/robustez_temporal_resumen.csv: rf, barajado, sin_macro_ni_month, auc, media
\newcommand{\aucRfSinMacroNiMonthBarajadoDesvio}{0{,}009}   % resultados/robustez_temporal_resumen.csv: rf, barajado, sin_macro_ni_month, auc, desvio
\newcommand{\recallRfAdelante}{0{,}243}                     % resultados/robustez_temporal_resumen.csv: rf, hacia_adelante, todas, recall_q, media
\newcommand{\caidaAucRfAdelante}{0{,}237}                   % resultados/robustez_temporal_resumen.csv: rf, barajado_menos_hacia_adelante, todas, auc, media
\newcommand{\deltaAucSinMacroRfBarajado}{\signoMenos 0{,}023} % resultados/robustez_temporal_resumen.csv: rf, barajado, sin_macro menos todas, auc, media
\newcommand{\deltaAucSinMacroNiMonthRfBarajado}{\signoMenos 0{,}059} % resultados/robustez_temporal_resumen.csv: rf, barajado, sin_macro_ni_month menos todas, auc, media
\newcommand{\aucKnnAdelante}{0{,}539}                       % resultados/robustez_temporal_resumen.csv: knn, hacia_adelante, todas, auc, media
\newcommand{\aucKnnAdelanteDesvio}{0{,}065}                 % resultados/robustez_temporal_resumen.csv: knn, hacia_adelante, todas, auc, desvio
\newcommand{\aucKnnSinMacroAdelante}{0{,}557}               % resultados/robustez_temporal_resumen.csv: knn, hacia_adelante, sin_macro, auc, media
\newcommand{\aucKnnSinMacroAdelanteDesvio}{0{,}099}         % resultados/robustez_temporal_resumen.csv: knn, hacia_adelante, sin_macro, auc, desvio
\newcommand{\aucKnnSinMacroBarajado}{0{,}754}               % resultados/robustez_temporal_resumen.csv: knn, barajado, sin_macro, auc, media
\newcommand{\aucKnnSinMacroBarajadoDesvio}{0{,}009}         % resultados/robustez_temporal_resumen.csv: knn, barajado, sin_macro, auc, desvio
\newcommand{\aucKnnSinMacroNiMonthAdelante}{0{,}542}        % resultados/robustez_temporal_resumen.csv: knn, hacia_adelante, sin_macro_ni_month, auc, media
\newcommand{\aucKnnSinMacroNiMonthAdelanteDesvio}{0{,}091}  % resultados/robustez_temporal_resumen.csv: knn, hacia_adelante, sin_macro_ni_month, auc, desvio
\newcommand{\aucKnnSinMacroNiMonthBarajado}{0{,}731}        % resultados/robustez_temporal_resumen.csv: knn, barajado, sin_macro_ni_month, auc, media
\newcommand{\aucKnnSinMacroNiMonthBarajadoDesvio}{0{,}009}  % resultados/robustez_temporal_resumen.csv: knn, barajado, sin_macro_ni_month, auc, desvio
\newcommand{\recallKnnAdelante}{0{,}215}                    % resultados/robustez_temporal_resumen.csv: knn, hacia_adelante, todas, recall_q, media
\newcommand{\caidaAucKnnAdelante}{0{,}245}                  % resultados/robustez_temporal_resumen.csv: knn, barajado_menos_hacia_adelante, todas, auc, media
\newcommand{\deltaAucSinMacroKnnBarajado}{\signoMenos 0{,}031} % resultados/robustez_temporal_resumen.csv: knn, barajado, sin_macro menos todas, auc, media
\newcommand{\deltaAucSinMacroNiMonthKnnBarajado}{\signoMenos 0{,}053} % resultados/robustez_temporal_resumen.csv: knn, barajado, sin_macro_ni_month menos todas, auc, media
\newcommand{\aucNbAdelante}{0{,}598}                        % resultados/robustez_temporal_resumen.csv: nb_categorico, hacia_adelante, todas, auc, media
\newcommand{\aucNbAdelanteDesvio}{0{,}087}                  % resultados/robustez_temporal_resumen.csv: nb_categorico, hacia_adelante, todas, auc, desvio
\newcommand{\aucNbSinMacroAdelante}{0{,}604}                % resultados/robustez_temporal_resumen.csv: nb_categorico, hacia_adelante, sin_macro, auc, media
\newcommand{\aucNbSinMacroAdelanteDesvio}{0{,}090}          % resultados/robustez_temporal_resumen.csv: nb_categorico, hacia_adelante, sin_macro, auc, desvio
\newcommand{\aucNbSinMacroBarajado}{0{,}749}                % resultados/robustez_temporal_resumen.csv: nb_categorico, barajado, sin_macro, auc, media
\newcommand{\aucNbSinMacroBarajadoDesvio}{0{,}012}          % resultados/robustez_temporal_resumen.csv: nb_categorico, barajado, sin_macro, auc, desvio
\newcommand{\aucNbSinMacroNiMonthAdelante}{0{,}566}         % resultados/robustez_temporal_resumen.csv: nb_categorico, hacia_adelante, sin_macro_ni_month, auc, media
\newcommand{\aucNbSinMacroNiMonthAdelanteDesvio}{0{,}101}   % resultados/robustez_temporal_resumen.csv: nb_categorico, hacia_adelante, sin_macro_ni_month, auc, desvio
\newcommand{\aucNbSinMacroNiMonthBarajado}{0{,}724}         % resultados/robustez_temporal_resumen.csv: nb_categorico, barajado, sin_macro_ni_month, auc, media
\newcommand{\aucNbSinMacroNiMonthBarajadoDesvio}{0{,}011}   % resultados/robustez_temporal_resumen.csv: nb_categorico, barajado, sin_macro_ni_month, auc, desvio
\newcommand{\recallNbAdelante}{0{,}287}                     % resultados/robustez_temporal_resumen.csv: nb_categorico, hacia_adelante, todas, recall_q, media
\newcommand{\caidaAucNbAdelante}{0{,}184}                   % resultados/robustez_temporal_resumen.csv: nb_categorico, barajado_menos_hacia_adelante, todas, auc, media
\newcommand{\deltaAucSinMacroNbBarajado}{\signoMenos 0{,}033} % resultados/robustez_temporal_resumen.csv: nb_categorico, barajado, sin_macro menos todas, auc, media
\newcommand{\deltaAucSinMacroNiMonthNbBarajado}{\signoMenos 0{,}058} % resultados/robustez_temporal_resumen.csv: nb_categorico, barajado, sin_macro_ni_month menos todas, auc, media
\newcommand{\aucSvmAdelante}{0{,}541}                       % resultados/robustez_temporal_resumen.csv: svm, hacia_adelante, todas, auc, media
\newcommand{\aucSvmAdelanteDesvio}{0{,}078}                 % resultados/robustez_temporal_resumen.csv: svm, hacia_adelante, todas, auc, desvio
\newcommand{\aucSvmSinMacroAdelante}{0{,}563}               % resultados/robustez_temporal_resumen.csv: svm, hacia_adelante, sin_macro, auc, media
\newcommand{\aucSvmSinMacroAdelanteDesvio}{0{,}122}         % resultados/robustez_temporal_resumen.csv: svm, hacia_adelante, sin_macro, auc, desvio
\newcommand{\aucSvmSinMacroBarajado}{0{,}754}               % resultados/robustez_temporal_resumen.csv: svm, barajado, sin_macro, auc, media
\newcommand{\aucSvmSinMacroBarajadoDesvio}{0{,}012}         % resultados/robustez_temporal_resumen.csv: svm, barajado, sin_macro, auc, desvio
\newcommand{\aucSvmSinMacroNiMonthAdelante}{0{,}549}        % resultados/robustez_temporal_resumen.csv: svm, hacia_adelante, sin_macro_ni_month, auc, media
\newcommand{\aucSvmSinMacroNiMonthAdelanteDesvio}{0{,}081}  % resultados/robustez_temporal_resumen.csv: svm, hacia_adelante, sin_macro_ni_month, auc, desvio
\newcommand{\aucSvmSinMacroNiMonthBarajado}{0{,}726}        % resultados/robustez_temporal_resumen.csv: svm, barajado, sin_macro_ni_month, auc, media
\newcommand{\aucSvmSinMacroNiMonthBarajadoDesvio}{0{,}009}  % resultados/robustez_temporal_resumen.csv: svm, barajado, sin_macro_ni_month, auc, desvio
\newcommand{\recallSvmAdelante}{0{,}233}                    % resultados/robustez_temporal_resumen.csv: svm, hacia_adelante, todas, recall_q, media
\newcommand{\caidaAucSvmAdelante}{0{,}233}                  % resultados/robustez_temporal_resumen.csv: svm, barajado_menos_hacia_adelante, todas, auc, media
\newcommand{\deltaAucSinMacroSvmBarajado}{\signoMenos 0{,}020} % resultados/robustez_temporal_resumen.csv: svm, barajado, sin_macro menos todas, auc, media
\newcommand{\deltaAucSinMacroNiMonthSvmBarajado}{\signoMenos 0{,}048} % resultados/robustez_temporal_resumen.csv: svm, barajado, sin_macro_ni_month menos todas, auc, media
\newcommand{\aucRfAdelanteFoldUno}{0{,}489}                 % resultados/robustez_temporal.csv: rf, hacia_adelante, todas, fold 1, validacion, auc
\newcommand{\pctYesBloqueUno}{4{,}7\,\%}                    % resultados/robustez_folds.csv: hacia_adelante, fold 1, pct_yes_validacion
\newcommand{\pctYesEntrenamientoAdelanteUno}{2{,}9\,\%}     % resultados/robustez_folds.csv: hacia_adelante, fold 1, pct_yes_train
\newcommand{\filasEntrenamientoAdelanteUno}{5\,490}         % resultados/robustez_folds.csv: hacia_adelante, fold 1, n_train
\newcommand{\aucRfAdelanteFoldDos}{0{,}500}                 % resultados/robustez_temporal.csv: rf, hacia_adelante, todas, fold 2, validacion, auc
\newcommand{\pctYesBloqueDos}{6{,}5\,\%}                    % resultados/robustez_folds.csv: hacia_adelante, fold 2, pct_yes_validacion
\newcommand{\pctYesEntrenamientoAdelanteDos}{3{,}8\,\%}     % resultados/robustez_folds.csv: hacia_adelante, fold 2, pct_yes_train
\newcommand{\filasEntrenamientoAdelanteDos}{10\,980}        % resultados/robustez_folds.csv: hacia_adelante, fold 2, n_train
\newcommand{\aucRfAdelanteFoldTres}{0{,}426}                % resultados/robustez_temporal.csv: rf, hacia_adelante, todas, fold 3, validacion, auc
\newcommand{\pctYesBloqueTres}{5{,}6\,\%}                   % resultados/robustez_folds.csv: hacia_adelante, fold 3, pct_yes_validacion
\newcommand{\pctYesEntrenamientoAdelanteTres}{4{,}7\,\%}    % resultados/robustez_folds.csv: hacia_adelante, fold 3, pct_yes_train
\newcommand{\filasEntrenamientoAdelanteTres}{16\,470}       % resultados/robustez_folds.csv: hacia_adelante, fold 3, n_train
\newcommand{\aucRfAdelanteFoldCuatro}{0{,}664}              % resultados/robustez_temporal.csv: rf, hacia_adelante, todas, fold 4, validacion, auc
\newcommand{\pctYesBloqueCuatro}{12{,}7\,\%}                % resultados/robustez_folds.csv: hacia_adelante, fold 4, pct_yes_validacion
\newcommand{\pctYesEntrenamientoAdelanteCuatro}{4{,}9\,\%}  % resultados/robustez_folds.csv: hacia_adelante, fold 4, pct_yes_train
\newcommand{\filasEntrenamientoAdelanteCuatro}{21\,960}     % resultados/robustez_folds.csv: hacia_adelante, fold 4, n_train
\newcommand{\aucRfAdelanteFoldCinco}{0{,}711}               % resultados/robustez_temporal.csv: rf, hacia_adelante, todas, fold 5, validacion, auc
\newcommand{\pctYesBloqueCinco}{35{,}2\,\%}                 % resultados/robustez_folds.csv: hacia_adelante, fold 5, pct_yes_validacion
\newcommand{\pctYesEntrenamientoAdelanteCinco}{6{,}5\,\%}   % resultados/robustez_folds.csv: hacia_adelante, fold 5, pct_yes_train
\newcommand{\filasEntrenamientoAdelanteCinco}{27\,450}      % resultados/robustez_folds.csv: hacia_adelante, fold 5, n_train
\newcommand{\filasBloqueAdelante}{5\,490}                   % resultados/robustez_folds.csv: hacia_adelante, fold 1, n_validacion
\newcommand{\pctParesEntreAnios}{66{,}3\,\%}                % resultados/oof_final_rf.csv con la y y el año inferido de train (cargar_train()): % de los pares «yes»–«no» que son de años distintos
\newcommand{\aucOofEntreAnios}{0{,}864}                     % resultados/oof_final_rf.csv con la y y el año inferido de train (cargar_train()): AUC fuera de fold de RF sobre los pares de años distintos
\newcommand{\aucOofDentroAnio}{0{,}658}                     % resultados/oof_final_rf.csv con la y y el año inferido de train (cargar_train()): AUC fuera de fold de RF sobre los pares del mismo año
\newcommand{\aucOofDosMilOcho}{0{,}598}                     % resultados/oof_final_rf.csv con la y y el año inferido de train (cargar_train()): AUC fuera de fold de RF en las filas de 2008
\newcommand{\aucOofDosMilNueve}{0{,}759}                    % resultados/oof_final_rf.csv con la y y el año inferido de train (cargar_train()): AUC fuera de fold de RF en las filas de 2009
\newcommand{\aucOofDosMilDiez}{0{,}749}                     % resultados/oof_final_rf.csv con la y y el año inferido de train (cargar_train()): AUC fuera de fold de RF en las filas de 2010
\newcommand{\recallRfDentroAnio}{0{,}378}                   % resultados/oof_final_rf.csv con la y y el año inferido de train (cargar_train()): RF llamando al 20 % de las filas de cada año por separado, «yes» alcanzados en los tres años / «yes» de train
\newcommand{\llamadasPorYesRfDentroAnio}{4{,}7}             % resultados/oof_final_rf.csv con la y y el año inferido de train (cargar_train()): RF llamando al 20 % de las filas de cada año por separado: llamadas / «yes» alcanzados
\newcommand{\recallRfDentroAnioDosMilOcho}{0{,}306}         % resultados/oof_final_rf.csv con la y y el año inferido de train (cargar_train()): RF llamando al 20 % de las filas de cada año por separado: recall en 2008
\newcommand{\llamadasDentroAnioDosMilOcho}{4\,423}          % resultados/oof_final_rf.csv con la y y el año inferido de train (cargar_train()): RF llamando al 20 % de las filas de cada año por separado: llamadas en 2008, llamadas(filas del año) de src/metricas.py
\newcommand{\yesDosMilOcho}{1\,089}                         % train (cargar_train()): filas con y = yes en 2008 (año inferido, como eda.html)
\newcommand{\recallRfDentroAnioDosMilNueve}{0{,}455}        % resultados/oof_final_rf.csv con la y y el año inferido de train (cargar_train()): RF llamando al 20 % de las filas de cada año por separado: recall en 2009
\newcommand{\llamadasDentroAnioDosMilNueve}{1\,832}         % resultados/oof_final_rf.csv con la y y el año inferido de train (cargar_train()): RF llamando al 20 % de las filas de cada año por separado: llamadas en 2009, llamadas(filas del año) de src/metricas.py
\newcommand{\yesDosMilNueve}{1\,758}                        % train (cargar_train()): filas con y = yes en 2009 (año inferido, como eda.html)
\newcommand{\recallRfDentroAnioDosMilDiez}{0{,}311}         % resultados/oof_final_rf.csv con la y y el año inferido de train (cargar_train()): RF llamando al 20 % de las filas de cada año por separado: recall en 2010
\newcommand{\llamadasDentroAnioDosMilDiez}{333}             % resultados/oof_final_rf.csv con la y y el año inferido de train (cargar_train()): RF llamando al 20 % de las filas de cada año por separado: llamadas en 2010, llamadas(filas del año) de src/metricas.py
\newcommand{\yesDosMilDiez}{864}                            % train (cargar_train()): filas con y = yes en 2010 (año inferido, como eda.html)
\newcommand{\recallTopeDosMilDiez}{0{,}385}                 % resultados/oof_final_rf.csv con la y y el año inferido de train (cargar_train()): RF llamando al 20 % de las filas de cada año por separado: llamadas / «yes» de 2010, el recall máximo en ese año, que tiene más «yes» que llamadas
\newcommand{\recallRfDentroMes}{0{,}295}                    % resultados/oof_final_rf.csv con la y y el año inferido de train (cargar_train()): RF llamando al 20 % de cada año y mes por separado, «yes» alcanzados / «yes» de train
\newcommand{\aucOofDentroMes}{0{,}564}                      % resultados/oof_final_rf.csv con la y y el año inferido de train (cargar_train()): AUC fuera de fold de RF sobre los pares del mismo año y mes (ponderado por pares, como \aucOofDentroAnio)
\newcommand{\aucRfBarajadoBloqueUno}{0{,}536}               % resultados/oof_final_rf.csv con la y de train (cargar_train()), en las filas del bloque 1 hacia adelante (resultados/robustez_folds.csv): AUC fuera de fold de RF (barajado)
\newcommand{\caidaAucRfBloqueUno}{0{,}047}                  % resultados/oof_final_rf.csv con la y de train (cargar_train()), en las filas del bloque 1 hacia adelante (resultados/robustez_folds.csv): AUC barajado menos el hacia adelante (resultados/robustez_temporal.csv: rf, todas, fold 1)
\newcommand{\yesEntrenamientoAdelanteUno}{159}              % train (cargar_train()): «yes» con fila <= fila_max_train del fold 1 hacia adelante (resultados/robustez_folds.csv)
\newcommand{\aucRfBarajadoBloqueDos}{0{,}519}               % resultados/oof_final_rf.csv con la y de train (cargar_train()), en las filas del bloque 2 hacia adelante (resultados/robustez_folds.csv): AUC fuera de fold de RF (barajado)
\newcommand{\caidaAucRfBloqueDos}{0{,}019}                  % resultados/oof_final_rf.csv con la y de train (cargar_train()), en las filas del bloque 2 hacia adelante (resultados/robustez_folds.csv): AUC barajado menos el hacia adelante (resultados/robustez_temporal.csv: rf, todas, fold 2)
\newcommand{\yesEntrenamientoAdelanteDos}{419}              % train (cargar_train()): «yes» con fila <= fila_max_train del fold 2 hacia adelante (resultados/robustez_folds.csv)
\newcommand{\aucRfBarajadoBloqueTres}{0{,}584}              % resultados/oof_final_rf.csv con la y de train (cargar_train()), en las filas del bloque 3 hacia adelante (resultados/robustez_folds.csv): AUC fuera de fold de RF (barajado)
\newcommand{\caidaAucRfBloqueTres}{0{,}158}                 % resultados/oof_final_rf.csv con la y de train (cargar_train()), en las filas del bloque 3 hacia adelante (resultados/robustez_folds.csv): AUC barajado menos el hacia adelante (resultados/robustez_temporal.csv: rf, todas, fold 3)
\newcommand{\yesEntrenamientoAdelanteTres}{778}             % train (cargar_train()): «yes» con fila <= fila_max_train del fold 3 hacia adelante (resultados/robustez_folds.csv)
\newcommand{\aucRfBarajadoBloqueCuatro}{0{,}748}            % resultados/oof_final_rf.csv con la y de train (cargar_train()), en las filas del bloque 4 hacia adelante (resultados/robustez_folds.csv): AUC fuera de fold de RF (barajado)
\newcommand{\caidaAucRfBloqueCuatro}{0{,}084}               % resultados/oof_final_rf.csv con la y de train (cargar_train()), en las filas del bloque 4 hacia adelante (resultados/robustez_folds.csv): AUC barajado menos el hacia adelante (resultados/robustez_temporal.csv: rf, todas, fold 4)
\newcommand{\yesEntrenamientoAdelanteCuatro}{1\,084}        % train (cargar_train()): «yes» con fila <= fila_max_train del fold 4 hacia adelante (resultados/robustez_folds.csv)
\newcommand{\aucRfBarajadoBloqueCinco}{0{,}762}             % resultados/oof_final_rf.csv con la y de train (cargar_train()), en las filas del bloque 5 hacia adelante (resultados/robustez_folds.csv): AUC fuera de fold de RF (barajado)
\newcommand{\caidaAucRfBloqueCinco}{0{,}051}                % resultados/oof_final_rf.csv con la y de train (cargar_train()), en las filas del bloque 5 hacia adelante (resultados/robustez_folds.csv): AUC barajado menos el hacia adelante (resultados/robustez_temporal.csv: rf, todas, fold 5)
\newcommand{\yesEntrenamientoAdelanteCinco}{1\,780}         % train (cargar_train()): «yes» con fila <= fila_max_train del fold 5 hacia adelante (resultados/robustez_folds.csv)

% --- Reserva · Curva de ganancia (H3) --------------------------------------------------------------
\newcommand{\gananciaRfDiez}{44{,}6\,\%}                    % resultados/ganancia_final.csv: rf, p = 10, pct_yes_alcanzado
\newcommand{\gananciaRfVeinte}{62{,}8\,\%}                  % resultados/ganancia_final.csv: rf, p = 20, pct_yes_alcanzado
\newcommand{\gananciaRfTreinta}{70{,}4\,\%}                 % resultados/ganancia_final.csv: rf, p = 30, pct_yes_alcanzado
\newcommand{\gananciaRfCincuenta}{81{,}7\,\%}               % resultados/ganancia_final.csv: rf, p = 50, pct_yes_alcanzado
\newcommand{\gananciaKnnVeinte}{61{,}7\,\%}                 % resultados/ganancia_final.csv: knn, p = 20, pct_yes_alcanzado
\newcommand{\gananciaNbVeinte}{61{,}4\,\%}                  % resultados/ganancia_final.csv: nb_categorico, p = 20, pct_yes_alcanzado
\newcommand{\gananciaSvmVeinte}{61{,}3\,\%}                 % resultados/ganancia_final.csv: svm, p = 20, pct_yes_alcanzado
\newcommand{\gananciaSinModeloVeinte}{20{,}0\,\%}           % resultados/ganancia_final.csv: sin_modelo, p = 20, pct_yes_alcanzado
\newcommand{\pctListaYesMitadRf}{12\,\%}                    % resultados/ganancia_final.csv: rf, el menor p con pct_yes_alcanzado >= 50
\newcommand{\pctListaYesOchentaRf}{48\,\%}                  % resultados/ganancia_final.csv: rf, el menor p con pct_yes_alcanzado >= 80

% GENERADO POR src/evaluar_test.py; no editar a mano.
% Marcadores: el test todavía no se evaluó (N0-1). La evaluación única
% los reemplaza: python3 -m src.evaluar_test
% \recalltest y \precisiontest reparten en proporción el empate en el corte, como
% src/metricas.py. \recallmatriztest y \precisionmatriztest salen de la matriz (\vptest,
% \fptest, \fntest), que lo desempata por fila. Difieren sólo si \llamadosempatadostest
% es menor que \empatadostest.
\newif\iftestpendiente\testpendientetrue
\newcommand{\auctest}{?}
\newcommand{\aucic}{?}
\newcommand{\recalltest}{?}
\newcommand{\recallic}{?}
\newcommand{\precisiontest}{?}
\newcommand{\fitest}{?}
\newcommand{\aptest}{?}
\newcommand{\vptest}{?}
\newcommand{\fptest}{?}
\newcommand{\fntest}{?}
\newcommand{\vntest}{?}
\newcommand{\recallmatriztest}{?}
\newcommand{\precisionmatriztest}{?}
\newcommand{\empatadostest}{?}
\newcommand{\llamadosempatadostest}{?}
\newcommand{\ntest}{?}
\newcommand{\llamadastest}{?}
\newcommand{\nevaluaciones}{?}


\begin{document}

% =====================================================================
% 1 · Portada. Sin \frametitle; el pie le pone el número de todos modos. Armada a mano, como en
% el TP1: volanta, título en bandera, filete azul y los datos en tres escalones de gris.
% =====================================================================
\begin{frame}
  \vspace{0.4em}
  {\fontsize{7.5}{9}\selectfont\volanta{\insertsubtitle}\par}
  \vskip1.0em
  {\usebeamerfont{title}\color{tinta}%
    \parbox[t]{0.80\textwidth}{\raggedright\inserttitle}\par}
  \vskip0.9em
  {\color{azul}\rule{4.5em}{2.2pt}\par}
  \vskip1.2em
  {\small\color{tintasec}\insertauthor\par}
  \vskip0.15em
  {\small\color{tintasec}\insertinstitute\par}
  \vskip1.5em
  {\footnotesize\color{apagado}%
    Dataset: \textit{Bank Marketing} --- \filasCsv{} contactos de campañas telefónicas de un
    banco, de \anioInicio{} a \anioFin\par}
  \vskip0.2em
  {\footnotesize\color{apagado}\insertdate\par}
  \note{Buenas tardes. Grupo siete: Trabajo Práctico 2, clasificación supervisada sobre
    \textit{Bank Marketing}.\\[0.4em]
    \textbf{Si preguntan.} Es el dataset que indica el enunciado: el \textit{Bank Marketing}, con
    campañas de \anioInicio{} a \anioFin. Las \filasCsv{} filas de la portada son el archivo entero,
    antes de quitar los \duplicados{} duplicados exactos (slide 3); desde la slide 2, cada número
    dice de qué conjunto sale.}
\end{frame}

% =====================================================================
\section{El problema}
% =====================================================================
% 2 · El problema. El diagrama entra en el orden del relato: qué se predice, con qué datos
% (lo que el banco sabe antes de marcar, D-07) y el desbalance, que es el título. La barra de
% clases está dibujada a escala con el % de «yes» de train (\pctYesTrain): si cambiara la
% partición, hay que mover el corte (hoy en 11,414 cm: el 88,7 % de los 4,3 cm de la barra).
\begin{frame}{El problema y el desbalance de clases}
  \vspace{0.4em}
  \centrado{%
  \begin{tikzpicture}[font=\small, x=1cm, y=1cm]
    % Las cuatro cajas, del mismo ancho y alto (text width fija el ancho aunque el renglón gris
    % sea más largo), a 1,2 cm de centro a centro: 0,22 cm de aire entre una y otra.
    \begin{scope}[every node/.append style={caja, text width=4.9cm, minimum height=0.98cm,
                                            visible on=<2->}]
      \node (cli) at (0,1.8)
        {el cliente\\[-2pt]{\scriptsize\gris{edad, trabajo, estudios, crédito}}};
      \node (his) at (0,0.6)
        {su historial\\[-2pt]{\scriptsize\gris{campañas anteriores y su resultado}}};
      \node (ctx) at (0,-0.6)
        {el contexto económico\\[-2pt]{\scriptsize\gris{euribor, empleo, precios, confianza}}};
      \node (lla) at (0,-1.8)
        {la llamada que se va a hacer\\[-2pt]{\scriptsize\gris{canal, mes, día, n.º de llamada}}};
    \end{scope}
    \node[cajaazul, minimum width=2.7cm, minimum height=1.2cm] (m) at (5.1,0)
      {modelo\\[-2pt]{\scriptsize\gris{antes de llamar}}};
    \node[cajanaranja, minimum width=4.3cm, minimum height=1.2cm] (y) at (9.75,0)
      {¿va a contratar?\\[-2pt]{\scriptsize\gris{\texttt{y}: yes o no}}};
    \draw[flecha] (m) -- (y);
    \draw[flecha, visible on=<2->] (cli.east) -- (m.160);
    \draw[flecha, visible on=<2->] (his.east) -- (m.173);
    \draw[flecha, visible on=<2->] (ctx.east) -- (m.187);
    \draw[flecha, visible on=<2->] (lla.east) -- (m.200);
    % La barra de clases, a escala: 4,3 cm = las \filasTrain{} filas de train; el corte, en el
    % 88,7 % (11,414 cm). El rótulo de arriba dice qué se cuenta; cada tramo lleva su clase.
    \begin{scope}[visible on=<3->]
      \node[anchor=south west, font=\footnotesize, text=tinta, inner sep=1pt]
        at (7.6,-1.12) {\textbf{respuestas en train}\sepdato{\color{tintasec}\filasTrain}};
      \fill[azul] (7.6,-1.72) rectangle (11.414,-1.17);
      \fill[naranja] (11.454,-1.72) rectangle (11.9,-1.17);
      \node[anchor=west, font=\footnotesize\bfseries, text=white, inner sep=1pt] at (7.72,-1.445)
        {«no»\sepdato\exactitudSiempreNo};
      \draw[naranja, line width=0.7pt] (11.677,-1.76) -- (11.677,-1.94);
      \node[anchor=north east, font=\footnotesize\bfseries, text=naranja, inner sep=1pt]
        at (11.9,-1.96) {«yes»\sepdato\pctYesTrain};
    \end{scope}
  \end{tikzpicture}}
  \note{Un banco ofrece un plazo fijo por teléfono. La pregunta es: antes de marcar, ¿este
    cliente va a contratar? Si lo sabemos, llamamos primero a los más probables.
    \avanza El modelo sólo puede usar lo que el banco sabe antes de la llamada: el cliente, su
    historial en campañas anteriores, el contexto económico y los datos de la llamada que va a
    hacer.
    \avanza Y la clase está muy desbalanceada: en train, sólo el \pctYesTrain{} dijo que sí.
    Decir siempre «no» ya acierta el \exactitudSiempreNo: la exactitud no sirve aquí.\\[0.4em]
    \textbf{Si preguntan.} ¿Cuántos son? \filasTrain{} filas de train (slide 3): \pctYesTrain{}
    «yes» y \exactitudSiempreNo{} «no». ¿Qué variable no entra? \texttt{duration}: se conoce al
    colgar (slide 4; D-05). ¿\texttt{campaign} existe antes de llamar? Sí: es el «n.º de llamada»
    de la llamada que se va a hacer, contando esa. La reserva es una inferencia: si cada fila
    guarda el total de la campaña y, tras un «yes», no se vuelve a llamar, su valor final depende
    en parte del resultado; el efecto no se midió (D-07). ¿Y \texttt{month},
    \texttt{day\_of\_week} y \texttt{contact}? Son el mes, el día y el canal de la llamada que se
    va a hacer, y se eligen antes de marcar (inferencia).}
\end{frame}

% =====================================================================
\section{Datos}
% =====================================================================
% 3 · Qué datos deciden (consejo «Importante» del enunciado, p. 3; guía A3): antes de la primera
% métrica, la partición y la tabla etapa → datos → qué decide. La barra de train y test está a
% escala del 80/20; las filas de la tabla entran de a una, y la del test va en naranja porque es
% la única que no decide nada.
\begin{frame}{División de los datos en train y test}
  \vspace{0.2em}
  \centrado{%
  \begin{tikzpicture}[x=1cm, y=1cm]
    \node[anchor=south west, font=\footnotesize, text=tintasec, inner sep=1pt] at (0,1.08)
      {\filasSinDuplicados{} filas sin duplicados, estratificadas por \texttt{y}};
    \fill[azul!20, rounded corners=2pt] (0,0) rectangle (11.25,1.0);
    \node[font=\small, align=center] at (5.625,0.5)
      {\textbf{train}\sepdato\proporcionTrain};
    \fill[naranja!20, rounded corners=2pt] (11.39,0) rectangle (14.2,1.0);
    \node[font=\small, align=center] at (12.795,0.5)
      {\textbf{test}\sepdato\proporcionTest};
    % La tabla: etapa, datos y qué decide.
    \node[anchor=west, inner sep=0pt] at (0.25,-0.62) {\cab{Etapa}};
    \node[anchor=west, inner sep=0pt] at (4.55,-0.62) {\cab{Datos}};
    \node[anchor=west, inner sep=0pt] at (8.85,-0.62) {\cab{Qué decide}};
    \draw[apagado, line width=0.4pt] (0,-0.86) -- (14.2,-0.86);
    \begin{scope}[visible on=<2->]
      \fill[rejilla!35] (0,-1.02) rectangle (14.2,-1.62);
      \node[anchor=west, font=\small] at (0.25,-1.32) {análisis de datos};
      \node[anchor=west, font=\small\bfseries, text=azul] at (4.55,-1.32) {train};
      \node[anchor=west, font=\small] at (8.85,-1.32) {qué variables entran, y cómo};
    \end{scope}
    \begin{scope}[visible on=<3->]
      \fill[rejilla!35] (0,-1.72) rectangle (14.2,-2.32);
      \node[anchor=west, font=\small] at (0.25,-2.02) {validación cruzada};
      \node[anchor=west, font=\small] at (4.55,-2.02) {\textcolor{azul}{\textbf{train}}, en 5 folds};
      \node[anchor=west, font=\small] at (8.85,-2.02) {modelo e hiperparámetros};
    \end{scope}
    \begin{scope}[visible on=<4->]
      \fill[naranja!12] (0,-2.42) rectangle (14.2,-3.02);
      \node[anchor=west, font=\small] at (0.25,-2.72) {evaluación final};
      \node[anchor=west, font=\small] at (4.55,-2.72)
        {\textcolor{naranja}{\textbf{test}}, una sola vez};
      \node[anchor=west, font=\small\bfseries, text=naranja] at (8.85,-2.72)
        {nada: sólo estima el desempeño};
    \end{scope}
  \end{tikzpicture}}
  \note{Primero, quitamos \duplicados{} duplicados exactos: si una copia cae en train y otra en
    test, el test deja de ser independiente. Después, 80/20 estratificado por la clase, con
    semilla fija.
    \avanza El análisis de datos, incluidas las pruebas con y sin cada variable, lee sólo train.
    \avanza La validación cruzada de 5 folds, también: con ella elegimos modelo e
    hiperparámetros.
    \avanza Y el test no decide nada: se abre una vez, al final, para estimar el desempeño.
    Cada número que sigue dice de qué conjunto sale.\\[0.4em]
    \textbf{Si preguntan.} ¿Por qué estratificar, si no se vio en clase? Con \pctYesTrain{} de
    «yes» (slide 2), una partición sin estratificar puede dejar al test con otra proporción sólo
    por la semilla; estratificando por \texttt{y}, train y test mantienen la misma (D-02). ¿Cómo se
    garantiza que nada mire el test? Por diseño, y la tabla lo fija: lo que decide lee sólo train,
    y el test se usa una vez, al final (D-04); lo que se aprende de los datos se ajusta dentro de
    cada fold (slide 7). ¿Cómo es cada vuelta? Valida con uno de los 5 folds y ajusta con los otros
    cuatro.}
\end{frame}

% =====================================================================
\section{EDA}
% =====================================================================
% 4 a 7 · El análisis exploratorio, cada observación con su consecuencia escrita (guía D1, gate
% G12): cada slide termina en un veredicto que nombra un efecto, no un método.
%
% 4 · La fuga. La figura es la ablación A1 (ΔAUC con duration) con la configuración de
% referencia, contra «sin modelo»; el veredicto es D-05. Es la primera slide con una métrica, y
% la sección Métricas viene después (slide 8): por eso la nota nombra aquí, en una pasada, el AUC
% de validación y su «sin modelo», y la 8 empieza por la métrica que el EDA ya usó (guía A2).
\begin{frame}{La duración predice el «yes», pero no existe antes de llamar}
  \renewcommand{\figgrandealto}{0.68}
  \centrado{\figgrande[0.97]{duration-techo.png}}
  \vspace{0.2em}
  \onslide<2->{\veredicto{\texttt{duration}: fuera del modelo}}
  \note{Primera decisión del análisis exploratorio: la fuga. \texttt{duration} es la duración
    de la última llamada, y se conoce al colgar, cuando el resultado ya se sabe. En train, el
    decil más corto, hasta \durationDecilCorto{} segundos, tiene \pctYesDurationDecilCorto{} de
    «yes»; el más largo, \pctYesDurationDecilLargo. Medimos cuánto pesa con el AUC de
    validación, la métrica con la que se compara todo el trabajo: \aucAzar{} es
    \rotuloLineaBase. Con \texttt{duration}, los cuatro saltan; Random Forest pasa de
    \aucReferenciaAblacionRf{} a \aucRfConDuration. Es un techo que no existe antes de llamar.
    \avanza Consecuencia: \texttt{duration} queda fuera del modelo.\\[0.4em]
    \textbf{Si preguntan.} ¿Con qué modelos? Los cuatro de la figura, con y sin
    \texttt{duration}, en AUC de validación, media de 5 folds. ¿Por qué el AUC antes de la slide de
    métricas? Porque ya es la métrica que decide (D-19): aquí se nombra, y en la slide 8 se
    justifica junto con el recall. ¿Qué otras variables se revisaron? Las de las cuatro cajas de
    la slide 2, que el banco conoce antes de marcar; \texttt{campaign}, con la reserva de esa slide
    (D-07). ¿Por qué medirla si sale? Porque el enunciado sugiere comparar con
    y sin las variables problemáticas, y el techo dice cuánto cuesta no tenerla.}
\end{frame}

% ---------------------------------------------------------------------
% 5 · El orden temporal (D-03, D-06). La figura es la del EDA en HTML, regenerada para
% proyectar: tasa de «yes» por bloque de filas del CSV original y el euribor debajo.
\begin{frame}{El archivo va por fecha: el «yes» sube de \pctYesPrimerBloqueEda{} a
  \pctYesUltimoBloqueEda}
  \renewcommand{\figgrandealto}{0.68}
  \centrado{\figgrande[0.97]{orden-temporal.png}}
  \vspace{0.2em}
  \onslide<2->{\veredicto{Partición aleatoria; la época va como limitación}}
  \note{Segunda: el archivo está ordenado por fecha, de mayo de 2008 a noviembre de 2010, y la
    tasa de «yes» no es estable: en bloques de \filasBloqueEda{} filas sube de
    \pctYesPrimerBloqueEda{} a \pctYesUltimoBloqueEda, mientras el euribor cae de
    \euriborPrimerBloqueEda{} a \euriborMinimoBloqueEda.
    \avanza Consecuencia: partimos al azar, como plantea el enunciado, y barajamos los folds,
    para que ninguno quede en una sola época. Lo que esto cuesta va como limitación, y es el
    hallazgo del final.\\[0.4em]
    \textbf{Si preguntan.} ¿Por año? La figura no lo da en números: en \anioInicio{} la tasa
    queda baja y pareja, y la subida es de 2009 y \anioFin, hasta \pctYesUltimoBloqueEda. ¿Por qué
    no partir por fecha? El enunciado plantea k-fold y no menciona el tiempo, y con el último
    \proporcionTest{} del archivo como test (slide 3), train quedaría casi entero en la época de
    tasa baja y test, en la de tasa alta (D-03). ¿Qué pasaba sin barajar? Los folds se armarían en
    el orden del archivo, y cada uno quedaría en una época; barajados, cada fold mezcla las épocas
    (slide 7; D-06). ¿De dónde sale el año? No viene en los datos, por eso la figura dice
    «inferido»; como el archivo va por fecha, cada vez que el mes retrocede empieza un año nuevo
    (inferencia desde el orden).}
\end{frame}

% ---------------------------------------------------------------------
% 6 · El centinela de pdays (D-10). El título lleva el 999 como macro (\pdaysCentinela), así pasa
% el verificador.
\begin{frame}{\pdaysCentinela{} no significa «nunca contactado»}
  \renewcommand{\figgrandealto}{0.68}
  \centrado{\figgrande[0.97]{pdays-999.png}}
  \vspace{0.2em}
  \onslide<2->{\veredicto{Decisión: \texttt{pdays} se mantiene sin transformar}}
  \note{Tercera: \texttt{pdays}. El \pdaysCentinela{} parece «no contactado antes», pero
    \filasCentinelaConPrevio{} filas con \pdaysCentinela{} sí tuvieron contactos previos, y los
    que tienen fecha contratan el \pctYesPdaysContactados.
    \avanza Se queda: escalado o en un árbol, el \pdaysCentinela{} funciona como un corte.\\[0.4em]
    \textbf{Si preguntan.} ¿Cuántos tienen días registrados? \filasPdaysContactados{} filas, con
    \pctYesPdaysContactados{} de «yes»; la barra del medio son las \filasCentinelaConPrevio{} con
    \pdaysCentinela{} y contactos previos, y su tasa se señala en pantalla. ¿Y una indicadora de
    contacto previo? Ya la da el historial (slide 2), y \texttt{previous} es la que separa los dos
    grupos del \pdaysCentinela; duplicarla pesaría doble en Naive Bayes y en las distancias de KNN
    (D-10). ¿Y los atípicos? Con casi todas las filas en \pdaysCentinela, el rango intercuartílico
    de \texttt{pdays} es cero: las \filasPdaysContactados{} filas con días registrados son justo
    sus atípicos, y no se borran (slide 7).}
\end{frame}

% ---------------------------------------------------------------------
% 7 · El pipeline: cuatro cajas iguales, cada una con lo que hace su paso. La partición queda
% afuera; lo que se ajusta con datos vive dentro de cada fold (enunciado, punto 1: «dentro del
% pipeline correspondiente»). Las cajas entran en el orden en que el pipeline las ejecuta, y el
% último overlay es el veredicto de la sección: los atípicos por IQR (D-17).
\begin{frame}{Todo se ajusta dentro de cada fold, en este orden}
  \vspace{0.3em}
  \centrado{%
  \begin{tikzpicture}[x=1cm, y=1cm, font=\small]
    % Cuatro cajas iguales; debajo del nombre, sólo lo que hace el paso.
    \begin{scope}[every node/.append style={text width=3.05cm, minimum height=1.2cm}]
      \node[caja, visible on=<1->] (a) at (1.7,0)
        {partición\\[-1pt]{\scriptsize\gris{5 folds barajados}}};
      \node[caja, visible on=<2->] (b) at (5.4,0)
        {reglas fijas\\[-1pt]{\scriptsize\gris{\columnasReferencia\salto\columnasFinales{}
         columnas}}};
      \node[cajaazul, visible on=<3->] (c) at (9.1,0)
        {codificar y escalar\\[-1pt]{\scriptsize\gris{con train del fold}}};
      \node[cajanaranja, visible on=<4->] (d) at (12.8,0)
        {clasificador\\[-1pt]{\scriptsize\gris{NB, SVM, KNN, RF}}};
    \end{scope}
    \draw[flecha, visible on=<2->] (a) -- (b);
    \draw[flecha, visible on=<3->] (b) -- (c);
    \draw[flecha, visible on=<4->] (c) -- (d);
    \begin{scope}[on background layer]
      \node[draw=apagado, dash pattern=on 2.4pt off 2pt, rounded corners=4pt, line width=0.6pt,
            inner sep=0.22cm, fit=(b) (c) (d), visible on=<2->] (fold) {};
    \end{scope}
    \node[anchor=south, font=\scriptsize, text=tintasec, inner sep=2pt, visible on=<2->]
      at (fold.north) {dentro de cada fold};
  \end{tikzpicture}}
  \vspace{0.9em}
  \onslide<5->{\veredicto{Atípicos por IQR: \filasAtipicasIqr{} filas marcadas, ninguna
    eliminada}}
  \note{Así queda el pipeline. La partición: \kFolds{} folds estratificados y barajados, así cada
    uno mezcla las épocas.
    \avanza Dentro de cada fold, reglas fijas: sale \texttt{duration}, \texttt{default} pasa a
    indicadora y se funden dos categorías raras; quedan \columnasFinales{} columnas.
    \avanza Después, la codificación de cada modelo: escalado para KNN y la SVM, deciles para
    Naive Bayes, ajustados sólo con el fold de entrenamiento.
    \avanza Los cuatro usan los mismos folds.
    \avanza Y los atípicos por IQR: se marcan \filasAtipicasIqr{} filas y no se borra ninguna,
    porque son señal.\\[0.4em]
    \textbf{Si preguntan.} ¿Qué hacen las reglas fijas? Son las mismas en todos los folds y no
    aprenden nada de los datos: sacan \texttt{duration} y funden categorías, de
    \columnasReferencia{} a \columnasFinales{} columnas. ¿Y Random Forest? Un árbol corta cada
    variable por separado: la escala no le cambia nada (D-12). ¿Cómo se decidió cada paso? Comparando
    en train (slide 3) el AUC con y sin cada decisión, como la de \texttt{duration} en la
    slide 4. ¿Por qué no borrar los atípicos? Entre los de
    \texttt{pdays}, el «yes» es \pctYesAtipicosPdays{} (slide 6), contra \pctYesTrain{} en todo
    train (slide 2): son los que más contratan (D-17). ¿Y KNN sin escalar? En train, \texttt{pdays}
    sería el \pctVarianzaPdays{} de la varianza (slide 16): la distancia miraría casi sólo esa
    variable.}
\end{frame}

% =====================================================================
\section{Métricas}
% =====================================================================
% 8 · Las dos métricas (D-19, D-20), cada una con su «sin modelo» debajo. El veredicto dice por
% qué no se usa la exactitud. El AUC ya se usó en el EDA (en la figura de la slide 4 y en lo que
% se dice en la 6), así que la nota no lo presenta como nuevo: empieza por la métrica que ya se
% vio y la justifica. La sección no se adelanta al EDA porque renumeraría las slides 4 a 8, a las
% que remiten el guion y el README de entrega.
\newcommand{\indicador}[4]{%
  \centrado{%
    \cab{#1}\\[0.5em]
    {\color{apagado}\rule{2.2em}{1pt}}\\[0.6em]
    {\fontsize{25}{28}\selectfont #2}\\[0.55em]
    {\small\color{tinta}#3}\\[0.4em]
    {\footnotesize\color{apagado}#4}}}
\begin{frame}{Dos métricas que no premian decir siempre «no»}
  \vspace{0.8em}
  \begin{columns}[T,onlytextwidth]
    \begin{column}{0.48\textwidth}
      \indicador{para comparar y elegir}{\dato{AUC}}%
        {¿ordena bien a los clientes?}{\rotuloLineaBase: \aucSinModelo}
    \end{column}
    \begin{column}{0.48\textwidth}
      \onslide<2->{\indicador{para medir el uso real}{\dato{recall al \presupuesto}}%
        {¿qué parte de los «yes» alcanzan las llamadas?}{\rotuloLineaBase: \recallSinModelo}}
    \end{column}
  \end{columns}
  \vspace{1.5em}
  \onslide<3->{\veredicto{La exactitud no se usa: decir siempre «no» ya da
    \exactitudSiempreNo}}
  \note{Las dos métricas son de la Clase 4, y la primera ya la usamos en el análisis
    exploratorio. El AUC mide si el modelo ordena bien a los clientes, que es justo priorizar a
    quién llamar; no depende de un umbral ni premia decir siempre «no»: sin modelo vale
    \aucAzar. Con ella comparamos y elegimos.
    \avanza La segunda es el uso real: si el banco llama al 20~\% de la lista con mayor puntaje,
    ¿qué parte de los que contratarían alcanza? Cuenta los clientes que habrían contratado y no se llamaron; sin
    modelo vale \recallSinModelo.
    \avanza La exactitud, no: decir siempre «no» ya da \exactitudSiempreNo.\\[0.4em]
    \textbf{Si preguntan.} ¿Por qué no el umbral de la Clase 4, el punto de la ROC más cercano
    a (0,1) (slide 52)? Es simétrico entre los dos errores, y aquí no cuestan lo mismo; además,
    con un umbral por modelo los recalls no se comparan: con el \presupuesto, todos hacen las
    mismas llamadas (D-20). ¿Por qué el \presupuesto? Es un supuesto de negocio, declarado como
    limitación (slide 17). ¿Y $F_1$ o la precisión? Con el \presupuesto{} fijo, precisión, recall
    y $F_1$ crecen con los «yes» alcanzados: ordenan igual a los modelos. ¿Y el desbalance? No se
    remuestrea: estas métricas ya no premian el «no» (D-21).}
\end{frame}

% =====================================================================
\section{Modelos}
% =====================================================================
% 9 · Los cinco clasificadores con valores de referencia (Naive Bayes en sus dos variantes,
% D-18), en dos pasos: la escala entera contra «sin modelo», con la ventana del zoom sombreada, y
% el zoom a esa ventana, donde se leen las barras de ±1 desvío.
\begin{frame}{Sin ajustar, gana \modeloMejorReferencia: \aucNbReferencia{} contra
  \aucSinModelo{} \rotuloLineaBase}
  \centrado{%
    \only<1>{\figgrande[0.97]{modelos-referencia.png}}%
    \only<2->{\figgrande[0.97]{modelos-referencia-zoom.png}}}
  \note{Primero, los modelos con valores de referencia, sin ajustar, en los mismos 5 folds:
    todos quedan muy por encima de sin modelo.
    \avanza De cerca, con su desvío entre folds: el mejor es Naive Bayes, en su versión
    categórica: discretizar las numéricas en deciles le suma \deltaAucNbCategorico{} sobre el
    gaussiano, porque ninguna numérica es normal. Y tres sobreajustan: Random Forest sin límite
    de profundidad, la SVM con RBF y KNN con k = \vecinosKnnReferencia. Lo muestran, y lo
    corrigen, las curvas que siguen.\\[0.4em]
    \textbf{Si preguntan.} ¿Cuánto sobreajustan? La figura sólo muestra validación: la brecha
    con train se ve en las curvas de las slides 10--12. ¿Qué es «sin ajustar»? Cada modelo con
    valores de referencia, antes de elegir un hiperparámetro por modelo con esas curvas. ¿Por qué
    el categórico? Le gana al gaussiano, \aucNbReferencia{} contra \aucNbGaussianoVal, porque las
    numéricas están lejos de ser normales: \texttt{pdays} es casi binaria (slide 6; D-18). ¿Por
    qué Random Forest da \aucReferenciaAblacionRf{} en la slide 4 y \aucRfReferencia{} aquí? Son
    dos mediciones de validación distintas: la de la slide 4 es del análisis exploratorio, y esta
    viene después del pipeline de la slide 7. Que aquella sea anterior al pipeline completo es
    inferencia por el orden; qué cambió no está en las slides.}
\end{frame}

% =====================================================================
\section{Hiperparámetros}
% =====================================================================
% 10 · La curva de sobreajuste (punto 3.1), revelada en tres pasos como la curva en U del TP1:
% train, validación, y las bandas con las zonas y el valor elegido. Mismo encuadre en las tres.
% «Sin modelo» (AUC 0,5) queda debajo del eje; sin línea de referencia, para que la curva
% ocupe el frame (pedido de menos texto de apoyo, 05/10).
\begin{frame}{¿Desde qué profundidad sobreajusta Random Forest?}
  \centrado{%
    \only<1>{\figgrande[0.97]{curva-rf-1.png}}%
    \only<2>{\figgrande[0.97]{curva-rf-2.png}}%
    \only<3->{\figgrande[0.97]{curva-rf-3.png}}}
  \note{Esta es la curva para discutir sobreajuste: la profundidad máxima de los árboles.
    Train sube siempre, y sin límite llega a \aucTrainRfProfundidadSinLimite: los árboles
    memorizan.
    \avanza La validación, en cambio, sube hasta \profundidadRfMejor{} y después baja, hasta
    \aucRfProfundidadSinLimite{} sin límite: la brecha se abre a
    \brechaRfProfundidadSinLimite. Eso es sobreajuste. En el otro extremo, con profundidad
    \profundidadRfMinima, train y validación casi empatan, \aucTrainRfProfundidadDos{} y
    \aucRfProfundidadDos: el modelo es demasiado simple, subajuste.
    \avanza Elegimos \profundidadRf: queda dentro de un error estándar del máximo y es más
    simple. Validación \aucRfProfundidadElegida, brecha \brechaRfProfundidadElegida.\\[0.4em]
    \textbf{Si preguntan.} ¿Por qué no \profundidadRfMejor, si ahí la validación es un poco más
    alta? La regla de un error estándar no elige el máximo: toma el más simple de los que quedan
    dentro de un error estándar de él (D-22). ¿Y por qué no la profundidad anterior a
    \profundidadRf, que se ve casi igual? Si hubiera quedado dentro, la regla la habría elegido: la
    estrella dice que quedó fuera, por poco. La banda no lo muestra: es $\pm 1$ desvío, más ancha
    que el error estándar, que no está dibujado. ¿De dónde sale la regla? No es de la cátedra;
    atenúa el sobreajuste a la validación (Clase 3, slides 98--100) porque no elige el máximo. ¿Y
    más árboles? No suman complejidad: no cambian el sesgo y sólo bajan la varianza (Clase 7,
    slide 61).}
\end{frame}

% ---------------------------------------------------------------------
% 11 · KNN: el sobreajuste está del otro lado del eje (Clase 9, slide 34 para la ponderación).
\begin{frame}{En KNN sobreajusta el k pequeño: se elige k = \vecinosKnn}
  \centrado{\figgrande[0.97]{curva-knn.png}}
  \note{En KNN el eje va al revés. Con \vecinosKnnMinimo{} vecino, train da \aucTrainKnnUno{} y
    validación \aucKnnUno: cada cliente es su propio vecino, sobreajuste puro. Al sumar vecinos
    la brecha se cierra, y desde \vecinosKnnMinimoUnoEs{} todo queda dentro de un error estándar
    del mejor; tomamos el más suave, \vecinosKnn: \aucKnnOchocientosUno, con brecha \brechaKnn.
    Ponderar por distancia queda peor en todo el rango.\\[0.4em]
    \textbf{Si preguntan.} ¿Dónde está el máximo? La validación forma una meseta en los k
    grandes: desde \vecinosKnnMinimoUnoEs{} casi no cambia, y se toma el más suave. ¿No es enorme
    k = \vecinosKnn? Con \pctYesTrain{} de «yes» (slide 2), estimar bien una proporción baja pide
    muchos vecinos (inferencia). ¿Y por distancia? Cada vecino pesa la inversa de su distancia
    (Clase 9, slide 34). Es la línea punteada: en train queda en el tope del eje, porque cada punto
    de train se encuentra a sí mismo, y en validación queda por debajo de la uniforme en todo el
    rango. ¿Un k alto no favorece a la clase mayoritaria (Clase 9, slide 31)? Sólo si se vota
    por mayoría. Aquí no se vota: la proporción de vecinos «yes» es el puntaje con que se ordena
    (slide 8).}
\end{frame}

% ---------------------------------------------------------------------
% 12 · La SVM: la curva de C con kernel lineal y, en el segundo paso, la del RBF atenuada, que
% es la que sobreajusta.
\begin{frame}{La SVM queda lineal y muy regularizada: C = \cSvm}
  \centrado{%
    \only<1>{\figgrande[0.97]{curva-svm.png}}%
    \only<2->{\figgrande[0.97]{curva-svm-2.png}}}
  \note{En la SVM exploramos C y el kernel. Con kernel lineal, la curva de C es plana, sin sobreajuste, y tomamos \cSvm, el más
    regularizado.
    \avanza Con RBF, en cambio, un C grande sobreajusta: train sube casi a uno y la validación
    cae. Y ni en su mejor C el RBF supera al lineal: queda lineal.\\[0.4em]
    \textbf{Si preguntan.} ¿Qué es C? El costo de las violaciones del margen: un C mayor castiga
    más y regulariza menos, como dice el eje (Clase 8, slide 49). ¿Por qué \cSvm? Con kernel
    lineal la curva es plana, y entre los que empatan se toma el más regularizado, que es el menor
    C del eje. ¿Y el RBF? Un C grande sobreajusta: train sube casi a uno y la validación cae; en su
    mejor C da \aucSvmBalanceadaMejor{} de AUC de validación (slide 16), debajo del \aucSvmVal{} del
    lineal. ¿Y $\gamma$? La curva del RBF varía sólo C, con $\gamma$ fijo; ajustarlo es la grilla
    conjunta de C y $\gamma$ (Clase 8, slides 48--49).}
\end{frame}

% =====================================================================
\section{Modelo final}
% =====================================================================
% 13 · Los cuatro ajustados, en AUC y en recall, contra «sin modelo» (D-23); como la 9, primero la
% escala entera y después el zoom.
\begin{frame}{Ajustados, gana \modeloFinal{} por \margenAucFinal{} de AUC en validación}
  \centrado{%
    \only<1>{\figgrande[0.97]{modelos-final.png}}%
    \only<2->{\figgrande[0.97]{modelos-final-zoom.png}}}
  \note{Con los hiperparámetros elegidos, los cuatro quedan lejos de sin modelo: entre Random
    Forest y sin modelo hay \margenAucSinModeloFinal{} de AUC.
    \avanza De cerca gana Random Forest: \aucRfVal{} de AUC y \recallRfVal{} de recall llamando
    al 20~\%. Le siguen KNN, \aucKnnVal; Naive Bayes, \aucNbVal, y la SVM, \aucSvmVal. Ninguno
    queda dentro de un error estándar de Random Forest. Pero la escala importa: entre el mejor y
    el peor hay apenas \rangoAucFinal.\\[0.4em]
    \textbf{Si preguntan.} ¿Y si hubiera empate? No lo hubo: ninguno quedó dentro de un error
    estándar, y Random Forest es primero también en recall: \recallRfVal{} contra \recallKnnVal{}
    de KNN. ¿Sobreajusta? Brecha de \brechaRf{} (slide 10). ¿Cuánto sumó ajustar? Contra la
    slide 9, Random Forest pasa de \aucRfReferencia{} a \aucRfVal, KNN de \aucKnnReferencia{} a
    \aucKnnVal{} y la SVM de \aucSvmRbfReferencia{} a \aucSvmVal. ¿Y la precisión? Con el
    \presupuesto{} fijo ordena igual que el recall (slide 8).\iftestpendiente{} Su inversa se ve en
    la slide 15: \llamadasPorYesRf{} llamadas por cada «yes», contra \llamadasPorYesSinModelo{} sin
    modelo.\else{} La slide 15 muestra el recall de test.\fi}
\end{frame}

% ---------------------------------------------------------------------
% 14 · El número de test, solo (guía C2, C3): la única slide del deck que no compara nada. Con
% el test pendiente, el hueco dibujado como tal; evaluado, el AUC con su intervalo. El intervalo
% (\aucic) va en modo texto: su «--» es una raya sólo ahí. La comparación con validación se dice,
% no se proyecta (C6): está en la nota.
\begin{frame}{¿Qué AUC esperar en clientes nuevos de la misma época?}
  \vspace{0.8em}
  \iftestpendiente
    \heropendiente{AUC de test}{pendiente: se evalúa una vez, después del 30/09}
  \else
    \heroe{\auctest}{AUC de \textbf{test}, sobre \filasTest{} clientes nuevos \sepdato
      intervalo del \nivelIntervalo: \aucic}
  \fi
  \note{¿Qué esperar en clientes nuevos? Lo dice el test: Random Forest reentrenado con todo
    train, sobre las \filasTest{} filas de test, una sola vez.
    \iftestpendiente Todavía no lo abrimos; lo haremos después de las respuestas de la
    cátedra.\else Da \auctest{} de AUC, con un intervalo del \nivelIntervalo{} de \aucic.\fi{}
    Y atención a qué mide: el test sale de las mismas campañas que train, así que estima
    clientes nuevos de esa época, no una campaña futura. Eso es el hallazgo.\\[0.4em]
    \textbf{Si preguntan.} ¿Por qué una sola vez? Si el test se mira más de una vez y algo cambia
    en el medio, deja de ser independiente: pasa a decidir, y la slide 3 dice que no decide nada
    (N0-1). \iftestpendiente\else ¿Qué dice el intervalo? Cuánto variaría el AUC de test con otra
    muestra de clientes nuevos de la misma época.\fi{} ¿Y contra validación? Se dice, no se
    proyecta: se nombran los dos valores, el de test y el \aucRfVal{} de validación (slide 13).
    ¿Por qué no aparece sin modelo? Porque el test no compite con nada: sólo mide (slide 3).}
\end{frame}

% ---------------------------------------------------------------------
% 15 · Qué errores comete (punto 5). Mientras el test siga cerrado, la matriz es la de
% validación fuera de fold, rotulada como tal y con su «sin modelo» (G11); cuando se evalúe, la
% misma slide pasa sola a la matriz de test, sin línea de base (C2). Es lo más honesto de las
% dos opciones del plan: la matriz de test con «?» no diría nada, y el análisis de errores de
% conclusiones.md (§3) se hizo sobre validación.
%
% La matriz de validación corta en su 20 % la lista completa fuera de fold, con los 5 folds
% juntos (\llamadasTrain llamadas), como la curva de ganancia y como se corta el test; la de
% conclusiones.md, §3, usa el mismo corte y da las mismas cuatro celdas. Los números del guion
% son los de numeros.md.
%
% Los dos ítems van en dos líneas con el corte elegido: la cifra arriba y, debajo, su lectura (en
% gris, la de «sin modelo»). Con un separador en la misma línea, la columna cortaba el segundo
% justo antes del «sin modelo», y el primero dejaba «caro» solo en su línea.
\newcommand{\matrizVP}{\iftestpendiente\vpRfVal\else\vptest\fi}
\newcommand{\matrizFP}{\iftestpendiente\fpRfVal\else\fptest\fi}
\newcommand{\matrizFN}{\iftestpendiente\fnRfVal\else\fntest\fi}
\newcommand{\matrizVN}{\iftestpendiente\vnRfVal\else\vntest\fi}
\newcommand{\matrizLlamadas}{\iftestpendiente\llamadasTrain\else\llamadastest\fi}
\newcommand{\matrizConjunto}{\iftestpendiente validación (OOF): los 5 folds juntos\else test\fi}
\newcommand{\matrizLista}{\iftestpendiente la lista completa\else la lista\fi}
\newcommand{\lecturaCosto}{\iftestpendiente\dato{\llamadasPorYesRf} llamadas por cada «yes»\\
  {\small\color{apagado}\rotuloLineaBase: \llamadasPorYesSinModelo}\else recall de test:
  \dato{\recallmatriztest}\fi}
\begin{frame}{Matriz de confusión de \modeloFinal{} llamando al 20~\%}
  \vspace{0.3em}
  \begin{columns}[c,onlytextwidth]
    \begin{column}{0.56\textwidth}
      \centrado{%
      \begin{tikzpicture}[x=1cm, y=1cm,
          celda/.style={minimum width=2.95cm, minimum height=1.3cm, align=center,
                        rounded corners=2pt, inner sep=2pt}]
        \node[anchor=south, font=\footnotesize, text=tintasec, inner sep=1pt, align=center]
          at (3.05,1.95) {\textbf{\matrizConjunto}\\llamando al 20~\% de \matrizLista:
          \matrizLlamadas{} llamadas};
        \node[inner sep=0pt] at (1.5,1.6) {\cab{llamado}};
        \node[inner sep=0pt] at (4.6,1.6) {\cab{no llamado}};
        \node[anchor=east, font=\small] at (-0.1,0.7) {«yes»};
        \node[anchor=east, font=\small] at (-0.1,-0.7) {«no»};
        \node[celda, fill=azul!12] at (1.5,0.7)
          {{\Large\bfseries\color{azul}\matrizVP}\\[-1pt]{\scriptsize\gris{«yes» alcanzados}}};
        \node[celda, fill=naranja!16, draw=naranja, line width=0.9pt] at (4.6,0.7)
          {{\Large\bfseries\color{naranja}\matrizFN}\\[-1pt]{\scriptsize\color{naranja}«yes» perdidos}};
        \node[celda, fill=rejilla!70] at (1.5,-0.7)
          {{\Large\color{tinta}\matrizFP}\\[-1pt]{\scriptsize\gris{llamadas de más}}};
        \node[celda, fill=rejilla!30] at (4.6,-0.7)
          {{\Large\color{tinta}\matrizVN}\\[-1pt]{\scriptsize\gris{bien descartados}}};
      \end{tikzpicture}}
    \end{column}
    \begin{column}{0.40\textwidth}
      \begin{itemize}\setlength\itemsep{1.1em}
        \item<2-> \ojo{\matrizFN{} «yes» sin llamar}:\\ el error más costoso
        \item<3-> \lecturaCosto
      \end{itemize}
    \end{column}
  \end{columns}
  \note{¿Qué errores comete? \iftestpendiente Con el test cerrado, la matriz es la de
    validación: los puntajes fuera de fold, sobre la lista completa de validación (los 5 folds
    juntos), llamando al 20~\%.\else Es la matriz de test, llamando al 20~\% de la
    lista.\fi{} Un falso positivo es una llamada de más; un falso negativo, un cliente que
    habría contratado y no se llama: es el error más costoso.
    \avanza \iftestpendiente Quedan \fnRfVal{} «yes» sin llamar, sobre todo de 2008 y sin
    historial.\else Quedan \fntest{} «yes» de test sin llamar.\fi
    \avanza \iftestpendiente Aun así, cada «yes» cuesta \llamadasPorYesRf{} llamadas, contra
    \llamadasPorYesSinModelo{} sin modelo.\else El recall de test, llamando al 20~\%, es
    \recallmatriztest.\fi\\[0.4em]
    \textbf{Si preguntan.} \iftestpendiente ¿Qué es OOF? Cada fila puntuada por el modelo del
    fold en que no entrenó: toda train, sin fuga. ¿Por qué la lista completa y no fold por fold?
    Los puntajes fuera de fold de los 5 folds se juntan en una lista y se corta su \presupuesto,
    \llamadasTrain{} llamadas: el mismo corte que se aplica al test. ¿De dónde sale el costo por
    «yes»? Las llamadas sobre los «yes» alcanzados, (\vpRfVal{} + \fpRfVal) / \vpRfVal; sin
    modelo, la inversa de la proporción de «yes», \pctYesTrain{} (slide 2).\else ¿Qué matriz es?
    La de test, con el \presupuesto{} de su lista, \llamadastest{} llamadas; no se compara con sin
    modelo (slide 14).\fi{} ¿Quiénes son los perdidos? La slide no los desglosa; lo que se ve lo
    hace esperable (inferencia): en 2008 el modelo ordena poco aun barajado (slide 19), y el
    \pctPoutcomeInexistente{} de train no tiene campaña previa (slide 17).}
\end{frame}

% =====================================================================
\section{Conclusiones}
% =====================================================================
% 16 · Por qué rindió cada uno (punto 5): una ficha por modelo, en el orden del ranking, con su
% AUC de validación en el color de su figura, el supuesto del modelo y un efecto medido. Son los
% cuatro ejemplos del enunciado: correlación y no normalidad (NB), escala (KNN, SVM) y relaciones
% no lineales (RF, SVM). El título pregunta y ninguna ficha da una causa del orden, porque ninguna
% medición la sostiene: la U de la edad se mezcla con la época (conclusiones.md, §1), el RBF, que
% podría doblarla, no mejora al lineal (D-22), y no se midió qué pierde cada modelo sin la edad.
% Lo que es inferencia lo dice la nota.
\tikzset{tarjeta/.style={fill=rejilla!45, filete=#1, text width=6.35cm, minimum height=1.75cm,
  inner xsep=0.85em, inner ysep=0.5em, align=flush left, font=\normalsize}}
\begin{frame}{¿Por qué ganó \modeloFinal{} y quedó última la \modeloUltimoFinal?}
  \vspace{0.5em}
  \centrado{%
  \begin{tikzpicture}[x=1cm, y=1cm]
    \node[tarjeta=colorrf, visible on=<1->] at (0,1.0)
      {\textbf{Random Forest}\sepdato\textcolor{colorrf}{\textbf{\aucRfVal}}\\[3pt]
       {\small no supone escala ni forma; es el que más pierde sin macro ni \texttt{month}:
        \deltaAucSinMacroNiMonthRfBarajado}};
    \node[tarjeta=colorknn, visible on=<2->] at (7.45,1.0)
      {\textbf{KNN}\sepdato\textcolor{colorknn}{\textbf{\aucKnnVal}}\\[3pt]
       {\small mide distancias: sin escalar, \texttt{pdays} es el \pctVarianzaPdays{} de
        la varianza}};
    \node[tarjeta=colornb, visible on=<3->] at (0,-1.0)
      {\textbf{Naive Bayes}\sepdato\textcolor{colornb}{\textbf{\aucNbVal}}\\[3pt]
       {\small supone independencia: tres macro con $r$ de \correlacionMacroMinima{} a
        \correlacionMacroMaxima}};
    \node[tarjeta=colorsvm, visible on=<4->] at (7.45,-1.0)
      {\textbf{SVM}\sepdato\textcolor{colorsvm}{\textbf{\aucSvmVal}}\\[3pt]
       {\small quedó lineal: el RBF no la mejora (\aucSvmBalanceadaMejor{} en su mejor C)}};
  \end{tikzpicture}}
  \vspace{0.6em}
  \referencia{AUC de validación \sepdato \rotuloLineaBase: \aucSinModelo}
  \note{¿Por qué rindió cada uno? Con los supuestos que pide el enunciado, y con lo que medimos.
    Random Forest no supone escala ni forma, y es el que más pierde sin las variables que fechan
    las filas: sin macro ni \texttt{month}, \deltaAucSinMacroNiMonthRfBarajado. Que gane por
    reconocer la época es inferencia; es lo que retoma el hallazgo.
    \avanza KNN mide distancias y necesita escalar: sin escalar, \texttt{pdays} tiene el
    \pctVarianzaPdays{} de la varianza.
    \avanza Naive Bayes supone independencia, y el bloque macro la viola, con correlaciones de
    \correlacionMacroMinima{} a \correlacionMacroMaxima; como las numéricas no son normales, el
    categórico le gana al gaussiano.
    \avanza La SVM quedó lineal porque el RBF no la mejora: en su mejor C da
    \aucSvmBalanceadaMejor. Por qué queda última no lo medimos, y entre los cuatro hay sólo
    \rangoAucFinal.\\[0.4em]
    \textbf{Si preguntan.} ¿Qué variables usa Random Forest? No medimos importancias; sí cuánto
    pierde sin las variables de época: \deltaAucSinMacroNiMonthRfBarajado{} sin macro ni
    \texttt{month}, el que más pierde de los cuatro (D-25). ¿Por qué el bloque macro no arruina a
    Naive Bayes si viola la independencia? Contar tres veces casi la misma información lleva las
    probabilidades a los extremos, pero no necesariamente cambia el orden, que es lo que mide el AUC
    (inferencia). ¿Por qué queda última la SVM? No lo medimos; la distancia con los demás es
    pequeña: de \aucSvmVal{} a \aucRfVal.}
\end{frame}

% ---------------------------------------------------------------------
% 17 · Limitaciones y mejoras, antes del hallazgo (guía A4, A7), cada lista ordenada por
% impacto y de tres fichas. El título dice la primera, que es la que el hallazgo explica.
\begin{frame}{Limitaciones: el test mide estas campañas, no una futura}
  \vspace{0.3em}
  \begin{columns}[T,onlytextwidth]
    \begin{column}{0.48\textwidth}
      \cab{Limitaciones}\par\vspace{0.6em}
      \ficha{La época: el test es de las mismas campañas}\par\vspace{0.5em}
      \ficha{Llamar al \presupuesto{} es un supuesto de negocio}\par\vspace{0.5em}
      \ficha{Hiperparámetros elegidos con la misma validación}
    \end{column}
    \begin{column}{0.48\textwidth}
      \onslide<2->{%
      \cab{Qué mejoraríamos}\par\vspace{0.6em}
      \ficha{Validar hacia adelante en el tiempo}\par\vspace{0.5em}
      \ficha{Más historial: el \pctPoutcomeInexistente{} no tiene campaña previa}\par
      \vspace{0.5em}
      \ficha{Un costo por tipo de error, no un \presupuesto{} fijo}}
    \end{column}
  \end{columns}
  \note{Tres limitaciones, por impacto. La época: el test sale de las mismas campañas, así que
    vale para clientes nuevos de esas campañas, no para una futura. Llamar al 20~\% es un
    supuesto. Y
    elegimos los hiperparámetros con la misma validación que los reporta: el máximo sale algo
    optimista.
    \avanza Lo primero que haríamos es validar hacia adelante en el tiempo, que es como se
    usaría el modelo, y sumar historial: el \pctPoutcomeInexistente{} no tiene campaña
    previa.\\[0.4em]
    \textbf{Si preguntan.} ¿Cuánto optimismo? No lo medimos: haría falta una validación cruzada
    anidada (Clase 3, slides 98--100), y aquí ese papel lo cumple el test, que se abre una vez; la
    regla de un error estándar lo atenúa, porque no elige el máximo (slide 10; D-22). ¿Y dentro de
    una campaña? El \aucRfVal{} y el \recallRfVal{} de validación (slide 13) ordenan una lista que
    mezcla las épocas; en los pares del mismo año, el AUC baja a \aucOofDentroAnio{} (slide 18).
    ¿Y \texttt{duration}? Con ella, Random Forest llega a \aucRfConDuration{} de AUC de
    validación, pero no existe antes de llamar (slide 4). ¿Otras mejoras? Calibrar las
    probabilidades, para poder usar un costo por tipo de error: hoy las dos métricas sólo miden el
    orden (slide 8).}
\end{frame}

% =====================================================================
\section{Hallazgo}
% =====================================================================
% 18 y 19 · El hallazgo H1 (la 19, como pesas: barajado arriba, hacia adelante abajo), lo último que se muestra (guía A4, A6), en dos pasos: primero la
% caída de los cuatro modelos y por qué el barajado es alto, lectura (a); después, bloque a
% bloque, las lecturas (b) y (c) (resultados/conclusiones.md, §2; D-25). Si el reloj va atrasado
% en la slide 13, el segundo paso de la 19 es lo que se recorta, nunca las limitaciones.
\begin{frame}{Entrenado con el pasado, Random Forest cae de \aucRfVal{} a \aucRfAdelante}
  \renewcommand{\figgrandealto}{0.68}
  \centrado{\figgrande[0.97]{robustez-modelos.png}}
  \vspace{0.2em}
  \onslide<2->{\veredicto{Barajado, el AUC es \aucOofEntreAnios{} entre años distintos y
    \aucOofDentroAnio{} dentro del mismo año}}
  \note{El hallazgo. Validamos hacia adelante dentro de train: cada fold entrena con el pasado y
    valida con el bloque siguiente. Barajados, los cuatro están entre \aucSvmVal{} y \aucRfVal;
    hacia adelante, ninguno pasa de \aucNbAdelante, y Random Forest cae a \aucRfAdelante.
    \avanza ¿Por qué tanto? Parte del \aucRfVal{} es distinguir años: en los pares «yes»–«no»
    de años distintos, el AUC es \aucOofEntreAnios; en los del mismo año, \aucOofDentroAnio.\\[0.4em]
    \textbf{Si preguntan.} ¿Cómo se validó hacia adelante? Con train ordenado por fecha y
    partido en bloques: cada fold valida con un bloque y entrena con todos los anteriores; son los
    bloques de la slide 19, todos dentro de train (D-25). ¿Por qué la barra hacia adelante es tan
    ancha? Porque sirve en dos bloques y en tres no (slide 19): los de 2008 dan de
    \aucRfAdelanteFoldTres{} a \aucRfAdelanteFoldDos, y los de 2009--2010,
    \aucRfAdelanteFoldCuatro{} y \aucRfAdelanteFoldCinco. La tasa de «yes» de cada bloque no lo
    explica: el AUC no depende de ella (Clase 4, slides 45--51). ¿Qué dice el veredicto? El AUC es
    la proporción de pares «yes»--«no» bien ordenados; con los puntajes barajados de validación, es
    \aucOofEntreAnios{} en los pares de años distintos y \aucOofDentroAnio{} en los del mismo año.
    No es entrenar en un año y validar en otro: eso lo mide la validación hacia adelante.}
\end{frame}

% ---------------------------------------------------------------------
\begin{frame}{En 2008 no sirve; en 2009--2010 pierde, pero no cae al azar}
  \renewcommand{\figgrandealto}{0.90}
  \centrado{%
    \only<1>{\figgrande[0.97]{robustez-folds.png}}%
    \only<2->{\figgrande[0.97]{robustez-folds-2.png}}}
  \note{Bloque a bloque: en los tres bloques de 2008, hacia adelante da \aucRfAdelanteFoldUno{} y
    \aucRfAdelanteFoldDos, y en el tercero, \aucRfAdelanteFoldTres, el orden incluso se invierte.
    \avanza Pero barajado, en esas mismas filas, tampoco pasa de \aucRfBarajadoBloqueTres: en
    2008 el modelo no sirve ni mezclando épocas. En 2009 y 2010, en cambio, pierde
    \caidaAucRfBloqueCuatro{} y \caidaAucRfBloqueCinco: una caída real, no el azar. Por eso el
    test vale para estas campañas.\\[0.4em]
    \textbf{Si preguntan.} ¿\aucRfAdelanteFoldTres{} es azar? Queda debajo de \aucAzar: en ese
    bloque el orden sale invertido (Clase 4, slides 45--51); pero es un solo bloque y sin
    intervalo, un indicio y no una prueba. ¿Es porque las variables macro fechan las filas? En
    parte: si el modelo sólo reconociera la época, frente a una nueva ordenaría como el azar, y eso
    dan los dos primeros bloques; el tercero, invertido, sugiere que además la relación entre las
    variables y el «yes» cambia con el tiempo (inferencia), con la misma reserva. ¿Por qué no
    validar así desde el principio? El enunciado plantea k-fold; un esquema temporal habría
    entrenado con los bloques de 2008, de \pctYesBloqueUno{} a \pctYesBloqueDos{} de «yes», para
    predecir el último, con \pctYesBloqueCinco{} (D-25). Es la primera mejora (slide 17).}
\end{frame}

% =====================================================================
\section{Preguntas}
% =====================================================================
% 20 · La pantalla que queda puesta durante las preguntas: el mismo vocabulario que la portada.
% La sección «Preguntas» no está en el índice del encabezado (este frame no lo muestra); sólo
% pone su nombre en el pie.
\begin{frame}
  \vspace{0.4em}
  {\fontsize{7.5}{9}\selectfont\volanta{Gracias}\par}
  \vskip1.0em
  {\usebeamerfont{title}\color{tinta}¿Preguntas?\par}
  \vskip0.9em
  {\color{azul}\rule{4.5em}{2.2pt}\par}
  \vskip1.2em
  {\small\color{tintasec}\insertauthor\par}
  \vskip0.15em
  {\small\color{tintasec}\insertinstitute\par}
  \vskip1.5em
  {\footnotesize\color{apagado}\insertsubtitle{} --- \insertdate\par}
  \note{Gracias; quedamos a disposición para las preguntas.\\[0.4em]
    \textbf{Las preguntas se responden con las slides de la charla}: se vuelve a la que corresponde y se
    contesta con lo que muestra. No hay slides de respaldo.}
\end{frame}

\end{document}
